\documentclass[11pt]{amsart}

\usepackage[utf8]{inputenc}
\usepackage[T1]{fontenc}
\usepackage{lmodern}
\usepackage{amsmath,amssymb,amsthm,mathtools}
\usepackage[margin=1.1in]{geometry}
\usepackage{enumitem}
\usepackage{booktabs}
\usepackage{array}
\usepackage{xcolor}
\usepackage[colorlinks=true,linkcolor=blue!50!black,citecolor=green!40!black,urlcolor=blue!60!black]{hyperref}
\usepackage[capitalize,nameinlink]{cleveref}
\emergencystretch=3em

\newtheorem{theorem}{Theorem}[section]
\newtheorem{lemma}[theorem]{Lemma}
\newtheorem{proposition}[theorem]{Proposition}
\newtheorem{corollary}[theorem]{Corollary}
\newtheorem{conjecture}[theorem]{Conjecture}
\newtheorem{question}[theorem]{Question}
\newtheorem*{mainthm}{Theorem 1}
\newtheorem*{mainthmp}{Theorem 1$'$}
\theoremstyle{definition}
\newtheorem{definition}[theorem]{Definition}
\theoremstyle{remark}
\newtheorem{remark}[theorem]{Remark}

\newcommand{\R}{\mathbb{R}}
\newcommand{\Dt}{\Delta_2}
\newcommand{\mult}{\mu}
\newcommand{\eps}{\varepsilon}
\newcommand{\conv}{\operatorname{conv}}
\newcommand{\turn}{\operatorname{turn}}
\newcommand{\dist}{\operatorname{dist}}
\newcommand{\Cc}{\mathrm{C}}      % circle
\newcommand{\Dd}{\mathrm{B}}      % closed disc
\newcommand{\chk}[1]{\textsuperscript{[#1]}}
\definecolor{RED}{rgb}{0.8,0,0}\definecolor{red}{rgb}{0.8,0,0}
\newcommand{\placeholder}[1]{\textcolor{red}{\textbf{[#1]}}}

\title[Second-largest and minimum distance multiplicities]{Second-largest and minimum distance multiplicities in planar point sets: an upper bound of $15/11$, and related results on Erd\H{o}s Problem \#132}

\author{Wingate Jones}

\date{Draft of 24 September 2026}

\begin{document}

\begin{abstract}
For a finite set $X\subset\R^2$ let $\Dt$ and $\delta$ be the second-largest and the smallest distance determined by $X$, and let $\mult(d)$ be the number of pairs of $X$ at distance $d$. Clemen, Dumitrescu and Liu asked for
$L=\limsup_{n\to\infty}\max_{|X|=n}\min\{\mult(\Dt),\mult(\delta)\}/n$; the known bounds were $9/7\le L\le 3/2$, the upper bound coming from Vesztergombi's inequality $\mult(\Dt)\le 3n/2$. We prove that $\min\{\mult(\Dt),\mult(\delta)\}\le \tfrac{15}{11}n+C_0$ for an absolute constant $C_0$, so $L\le 15/11<3/2$; in one range of the parameters the bound improves to $\tfrac43n+C_0$. We first give a shorter proof of the weaker bound $\tfrac{54}{37}n+C_0$; that weaker bound is formalised in Lean~4.

We also give two results on Erd\H{o}s Problem \#132, which asks whether some distance other than the diameter must occur at most $n$ times. First, if no vertex of the convex hull has four points of $X$ at distance $\Dt$ (for instance, if no four points of $X$ are concyclic), then $\mult(\Dt)\le\tfrac32|L_1|+|L_2|$, where $L_1,L_2$ are the first two convex layers; so, under this hypothesis, the answer is yes whenever at least half as many points lie at depth $\ge 3$ as on the hull. This rests on a layer bound whose inputs from Vesztergombi's structure theory we prove directly. Second, for $n$ points in convex position there are at least three distances occurring between $1$ and $n$ times when $n$ is even and at least $6$, and when $n\in\{7,9,11,13\}$. The cases $n=11,13$ are computer-assisted: they reduce to the unsatisfiability of finite ordinal problems, certified by DRAT/LRAT proof checkers, and for $n=7,11,13$ also by a Lean~4 formalisation.
\end{abstract}

\maketitle

\section*{AI disclosure}
The author used an AI assistant (Claude, Anthropic) for exploration, computation, formalisation and drafting. The checks performed on the results are described in \cref{app:verification}.

\section{Introduction}\label{sec:intro}

Let $X\subset\R^2$ be a set of $n$ points. Write $\Delta=\Delta(X)$ for its diameter, $\Dt$ for its second-largest distance and $\delta$ for its minimum distance, and write $\mult(d)=\mult_X(d)$ for the number of unordered pairs of points of $X$ at distance $d$. Hopf and Pannwitz~\cite{HP34} proved that $\mult(\Delta)\le n$.

\subsection*{Erd\H{o}s Problem \#132}
Erd\H{o}s and Pach asked whether there must be a \emph{second} distance, besides the diameter, that occurs at least once and at most $n$ times, and whether the number of such distances must tend to infinity (perhaps being at least $n^{1-o(1)}$). This is Problem \#132 on the Erd\H{o}s problems website~\cite{EP132}, where it is recorded as open and where Erd\H{o}s's offer of \$100 for ``any nontrivial result'' is recorded (reference [Er97e] there). The statement fails for $n=4$ (two equilateral triangles glued along an edge) and holds for $n=5,6$ by Erd\H{o}s and Fishburn~\cite{EF95a}. We call a distance $d$ \emph{rare} if $1\le\mult(d)\le n$.

Clemen, Dumitrescu and Liu~\cite{CDL25} proved that the answer to the first question is yes for sets in convex position (their Theorem~1.2), and that $\mult(\Dt)\le n$ whenever
\begin{equation}\label{eq:cdl13}
\min\Bigl\{\tfrac32(|L_1|+|L_2|),\ \tfrac43|L_1|+2|L_2|,\ 2|L_1|+|L_2|\Bigr\}\le n,
\end{equation}
where $L_1,L_2$ are the first two convex layers (their Theorem~1.3). They showed that neither $\Dt$ nor $\delta$ is always rare: they constructed sets with $\min\{\mult(\Dt),\mult(\delta)\}\ge 9n/8-o(n)$ (their Proposition~1.5). They asked (their Problem~1.6) for the value of
\[
L:=\limsup_{n\to\infty}\ \max_{|X|=n}\ \frac{\min\{\mult(\Dt),\mult(\delta)\}}{n}.
\]
Vesztergombi~\cite{Ve87} proved that $\mult(\Dt)\le\frac32 n$ for every $n$-point planar set, which gives $L\le 3/2$. On the Erd\H{o}s problems forum, the user ienjoymath~\cite{ienjoymath} gave a construction showing $L\ge 9/7$ and conjectured $L=9/7$. As far as we know, $9/7\le L\le 3/2$ were the best bounds before this paper.

\subsection*{Results}
Our main result is an upper bound strictly below $3/2$.

\begin{mainthm}
There is an absolute constant $C_0$ such that every $n$-point set $X\subset\R^2$ satisfies
\[
\min\{\mult(\Dt),\mult(\delta)\}\ \le\ \tfrac{15}{11}\,n+C_0 .
\]
Consequently $L\le 15/11\approx 1.3636$. Moreover, put $t:=(\Delta-\Dt)/\delta$ and $\tau:=t-\delta(1-t^2)/(2\Dt)$. If $\Delta\le1.94\,\Dt$ and $\tau\ge\sqrt3/2+3/100$, then $\min\{\mult(\Dt),\mult(\delta)\}\le\frac43n+C_0$.
\end{mainthm}
(Here $X$ is assumed to determine at least two distances; otherwise $n\le3$.)

The constant $C_0$ is explicit; the proof gives $C_0=3N_0+2\cdot10^8$ with $N_0=1.51\cdot 10^9$ (\cref{sec:proof15}). We have made no attempt to optimise it. The range in the second statement is ``Region~II'' of \cref{sec:regions}; the quantity $\tau$ is the one defined in \cref{sec:setup} after scaling to $\delta=1$. When $\Delta>1.94\,\Dt$ we have $\mult(\Dt)\le n+1$ (\cref{prop:degenerate}).

We also prove the weaker bound
\begin{mainthmp}
There is an absolute constant $C_0$ such that every $n$-point set $X\subset\R^2$ satisfies $\min\{\mult(\Dt),\mult(\delta)\}\le \tfrac{54}{37}\,n+C_0$.
\end{mainthmp}
Theorem~1$'$ follows from Theorem~1, since $15/11<54/37\approx1.4595$. We give its proof first (\cref{sec:setup,sec:reductions,sec:bad,sec:regions,sec:degenerate,sec:assembly}) because it is shorter and self-contained, and because Theorem~1$'$ is the statement formalised in Lean (\cref{app:verification}). The proof of Theorem~1 (\cref{sec:upgrade,sec:loss1,sec:loss2,sec:proof15}) reuses the set-up, the bad-edge bound, Region~I, the mixed T-edge lemma and the degenerate regime, and replaces the treatment of Regions~II and~III (in particular the ownership/pairing argument) and the inequality~(E) of \cref{sec:reductions}. Theorem~1 is not formalised.

\Cref{thm:layer} and \cref{cor:A1} give a new class of sets for which the answer to the first question of Problem \#132 is yes. Let $m_3:=n-|L_1|-|L_2|$ be the number of points of depth at least $3$. Condition~\eqref{eq:cdl13} contains the case $2|L_1|+|L_2|\le n$, i.e.\ $|L_1|\le m_3$. We show that $|L_1|\le 2m_3$ suffices under a genericity hypothesis that holds, for instance, when no four points of $X$ are concyclic. The inputs we need from Vesztergombi's structure theory~\cite{Ve96} are proved here from scratch (\cref{prop:core2,thm:deg4}).

\Cref{sec:convex} treats the second question of Problem \#132 in convex position. \Cref{thm:convex-summary} states that a set of $n$ points in convex position has at least three rare distances for every $n$ with $6\le n\le 13$. Some of these cases follow from the literature, some follow from a parity-sharp version of the counting argument of~\cite{CDL25}, and the cases $n=11,13$ are new and computer-assisted. In \cref{sec:open} we list open problems.

\subsection*{Idea of the proof of Theorem 1}
Let $S$ be the set of points that have a partner at distance $\Dt$, let $R:=X\setminus S$, and let $Q\subseteq S$ be the set of points with exactly one partner at distance $\Dt$. Vesztergombi's bound applied to $S\setminus Q$ gives $\mult(\Dt)\le\frac32|S|-\frac12|Q|$. On the other side we use the exact identity $\mult(\delta)=e(S)+\sum_{y\in R}(\deg y+m_y)/2$, where $e(S)$ is the number of $\delta$-pairs inside $S$, $\deg y$ is the number of points at distance $\delta$ from $y$, and $m_y$ is the number of those that lie in $S$. Almost all points of $S$ lie on the boundary of a convex body: either the ``ball polygon'' $K=\bigcap_{w\in X}\Dd(w,\Dt)$, or the convex hull. Each such point has a well-defined inward direction. We show that all but $O(1)$ of the $\delta$-edges inside $S$ join points with almost equal inward directions. Such edges come in two kinds: tangential edges, and \emph{rungs} joining the two boundaries, whose geometry is fixed by an exact identity (\cref{lem:identities}(c)). We show that $e(S)\le|S|+|Q|+O(1)$: apart from $O(1)$ exceptions, a rung either ends at a point of $Q$ or forces a tangential edge to be missing. For the points of $R$ a local analysis of $\partial K$ shows $\deg y+m_y\le10$ for all but $O(1)$ of them ($\le8$ in Region~II). The weights $8/11$ and $3/11$ then give $15/11$ (and $2/3$, $1/3$ give $4/3$ in Region~II). The regime where the diameter is much larger than $\Dt$ is handled separately by an incidence argument.

The weaker Theorem~1$'$ uses only the crude bound $\mult(\delta)\le e(S)+6|R|$, linear combinations of it with $\mult(\Dt)\le\frac32|S|$ and a weighted variant (\cref{lem:lp,lem:lp2}), and an ownership/pairing argument that bounds the density of rungs, giving $e(S)\le\frac43|S|+O(1)$.

\medskip\noindent\textbf{Notation.} $\Dd(c,r)$ and $\Cc(c,r)$ are the closed disc and the circle with centre $c$ and radius $r$. For unit vectors $u,u'$, $\angle(u,u')\in[0,\pi]$ is the angle between them; recall $|u-u'|\le\angle(u,u')$. Angles are in radians unless marked with a degree sign.

\section{Set-up for the proof of Theorem 1}\label{sec:setup}

We scale so that $\delta=1$. Throughout \cref{sec:setup,sec:reductions,sec:bad,sec:regions,sec:degenerate,sec:assembly,sec:upgrade,sec:loss1,sec:loss2,sec:proof15}, $X$ is an $n$-point set with $n\ge 4$ and at least two distinct distances, so $\Delta>\Dt\ge\delta=1$. We put
\[
t:=\Delta-\Dt>0,\qquad \tau:=t-\frac{1-t^2}{2\Dt},\qquad \beta:=\frac1{100}.
\]
Then $\tau-t=(t^2-1)/(2\Dt)$. Hence $\tau\le t$ when $t\le 1$, and $\tau>1$ if and only if $t>1$.\chk{C2} A crucial feature of $X$ is that it has \emph{no distance in the open interval $(\Dt,\Delta)$}.

We use the following objects.
\begin{itemize}[leftmargin=2em]
\item $D$ is the set of endpoints of diametral pairs. Each point of $D$ is a vertex of $P:=\conv X$.
\item $S$ is the set of points of $X$ that have a partner at distance $\Dt$; $R:=X\setminus S$ and $r:=|R|$, $s:=|S|$.
\item $K:=\bigcap_{w\in X}\Dd(w,\Dt)$, a compact convex set (a ``ball polygon'').
\item $Q\subseteq S$ is the set of points of $X$ that have exactly one partner at distance $\Dt$ (their \emph{$\Dt$-degree} is $1$), and $p:=|Q|$.
\item $e(Y)$ is the number of $\delta$-pairs inside $Y\subseteq X$.
\end{itemize}
The \emph{degenerate regime} is $\Delta>1.94\,\Dt$ (\cref{sec:degenerate}); the \emph{nondegenerate regime} is $\Delta\le 1.94\,\Dt$ (\cref{sec:reductions,sec:bad,sec:regions}).

\begin{lemma}[ball polygon]\label{lem:ballpoly}
\begin{enumerate}[label=(\alph*)]
\item $X\setminus D\subseteq K$.
\item $S\setminus D\subseteq\partial K$.
\item $\dist(z,K)\ge t$ for every $z\in D$.
\end{enumerate}
\end{lemma}
\begin{proof}
(a) A point $y\notin D$ has all its distances at most $\Dt$, since there are none in $(\Dt,\Delta)$. (b) If $v\in S\setminus D$ and $|vw|=\Dt$, then $v\in K\cap \Cc(w,\Dt)\subseteq\partial K$. (c) If $z\in D$ and $|zx|=\Delta$, then $K\subseteq\Dd(x,\Dt)$, while $|zx|=\Dt+t$.
\end{proof}

\subsection*{Direction field}
For each $p\in S$ we fix a unit vector $u_p$.
\begin{itemize}[leftmargin=2em]
\item If $p\in S\setminus D$, choose a $\Dt$-partner $w_p$ and put $u_p:=(w_p-p)/\Dt$. Since $K\subseteq\Dd(w_p,\Dt)$, $u_p$ is an inward normal of $K$ at $p$.
\item If $p\in S\cap D$, choose a diametral partner $x_p$ and put $u_p:=(x_p-p)/\Delta$. This is an inward normal of $P$ at $p$.
\end{itemize}

\begin{lemma}[exact identities]\label{lem:identities}
Let $y\in X$, $y\neq p$.
\begin{enumerate}[label=(\alph*)]
\item If $p\in S\cap D$, then $(y-p)\cdot u_p\ge |y-p|^2/(2\Delta)>0$.
\item If $p\in S\setminus D$, then either $(y-p)\cdot u_p\ge |y-p|^2/(2\Dt)$ (we say $y$ is \emph{normal} for $p$), or $|y-w_p|=\Delta$, so that $y\in D$ (we say $y$ is an \emph{exception} for $p$).
\item In the exception case with $|y-p|=1$ we have exactly $(y-p)\cdot u_p=-\tau$. Hence $y-p=-\tau u_p\pm\sigma u_p^{\perp}$ with $\sigma:=\sqrt{1-\tau^2}$. In particular exceptions at distance $1$ exist only when $\tau\le 1$.
\end{enumerate}
\end{lemma}
\begin{proof}
(a) $|y-x_p|^2=|y-p|^2-2\Delta(y-p)\cdot u_p+\Delta^2\le\Delta^2$. (b) $|y-w_p|^2=|y-p|^2-2\Dt(y-p)\cdot u_p+\Dt^2$, and $|y-w_p|\in[0,\Dt]\cup\{\Delta\}$. (c) With $|y-p|=1$ and $|y-w_p|=\Delta=\Dt+t$ the same identity gives $(y-p)\cdot u_p=(1+\Dt^2-\Delta^2)/(2\Dt)=(1-2\Dt t-t^2)/(2\Dt)=-\tau$.\chk{C2}
\end{proof}

\begin{lemma}[no common normal]\label{lem:nocommon}
If $p\neq q$ both lie in $S\setminus D$, then $u_p\neq u_q$.
\end{lemma}
\begin{proof}
Neither point lies in $D$, so each is normal for the other: $(q-p)\cdot u_p>0$ and $(p-q)\cdot u_q>0$. These cannot both hold if $u_p=u_q$.
\end{proof}

\subsection*{Good edges, T-edges and rungs}
A $\delta$-edge $pq$ inside $S$ is \emph{good} if $\angle(u_p,u_q)\le\beta$, and \emph{bad} otherwise. For a good edge put $e:=q-p$, $a:=e\cdot u_p$ and $b:=-e\cdot u_q$. Then $|a+b|=|e\cdot(u_p-u_q)|\le\beta$.

\begin{lemma}[dichotomy]\label{lem:dichotomy}
Assume $\tau>2\beta$. Each good edge is of exactly one of the following two kinds.
\begin{itemize}[leftmargin=2em]
\item \emph{T-edge:} $a,b\in(0,\beta)$.
\item \emph{Rung:} one end $v\in S\setminus D$ sees the other end $z\in S\cap D$ as an exception. Then $(z-v)\cdot u_v=-\tau$ and $(v-z)\cdot u_z\in[\tau-\beta,\tau+\beta]$.
\end{itemize}
\end{lemma}
\begin{proof}
By \cref{lem:identities}, from each end the other end is either normal (then the corresponding one of $a,b$ is positive) or an exception (then it equals $-\tau$). An exception target lies in $D$, and points of $D$ see everything as normal, so both ends cannot be exceptions. If neither is, then $a,b>0$ and $a+b\le\beta$. If one is, say $a=-\tau$, then $b\in[\tau-\beta,\tau+\beta]$, and $b>0$ shows that $b$ is not an exception value.
\end{proof}

We say a T-neighbour $q$ of $p$ is on the \emph{right} (resp.\ \emph{left}) of $p$ if $(q-p)\cdot u_p^{\perp}>0$ (resp.\ $<0$), where $u^\perp$ denotes $u$ rotated by $-90^\circ$.
\begin{itemize}[leftmargin=2em]
\item[(T1)] Each $p$ has at most one T-neighbour on each side. Indeed, two T-directions on one side both make an angle in $(0,\arcsin\beta)$ with $\pm u_p^\perp$, so they are less than $60^\circ$ apart, whereas two points at distance $1$ from $p$ and at distance $\ge 1$ from each other are at least $60^\circ$ apart as seen from $p$.
\item[(T2)] If $q$ is the right T-neighbour of $p$ then $p$ is the left T-neighbour of $q$, since $\angle(u_p,u_q)\le\beta$.
\end{itemize}
Hence the T-edges form vertex-disjoint directed paths and cycles, and every cycle has length at least $3$ (a vertex cannot be both the left and the right T-neighbour of another).

\subsection*{Turning}
Let $\gamma$ be a closed arc of the boundary of a compact convex set with nonempty interior. Its \emph{turning} $\turn(\gamma)$ is the length of the set $N(\gamma)\subseteq S^1$, the union of the normal cones at \emph{all} points of $\gamma$, including the full normal cones at its endpoints.

\begin{lemma}[multiplicity]\label{lem:multiplicity}
Let $\gamma_1,\dots,\gamma_m$ be closed boundary arcs of one convex body, and suppose every boundary point lies in at most $M$ of them. Then $\sum_i\turn(\gamma_i)\le 2\pi M$.
\end{lemma}
\begin{proof}
For all but countably many directions $\nu\in S^1$ the set of boundary points with outer (or inner) normal $\nu$ is a single point $y_\nu$, and $\nu\in N(\gamma_i)$ if and only if $y_\nu\in\gamma_i$. Integrate over $\nu$.
\end{proof}

Points of $S\setminus D$ lie on $\partial K$ and points of $S\cap D$ lie on $\partial P$; in each case distinct points are at Euclidean, hence arc-length, distance at least $1$. So a boundary arc of length $\ell$ contains at most $\ell+1$ of them, and a boundary of perimeter $\ell$ contains at most $\ell$ of them.

\section{Reductions}\label{sec:reductions}

\begin{lemma}[(VS)]\label{lem:VS}
For every $S'\subseteq S$, the number of $\Dt$-pairs inside $S'$ is at most $\frac32|S'|$. In particular $\mult(\Dt)\le\frac32 s$.
\end{lemma}
\begin{proof}
Every $\Dt$-pair of $X$ lies in $S$, and every distance within $S'$ is a distance of $X$. If $S'$ contains no $\Dt$-pair there is nothing to prove. If $S'$ contains a diametral pair of $X$, then $\operatorname{diam}S'=\Delta$ and, because $X$ has no distances in $(\Dt,\Delta)$, $\Dt$ is the second-largest distance of $S'$; Vesztergombi's theorem~\cite{Ve87} gives at most $\frac32|S'|$ such pairs. Otherwise every distance in $S'$ is at most $\Dt$, so $\Dt=\operatorname{diam}S'$ and Hopf--Pannwitz~\cite{HP34} gives at most $|S'|$.
\end{proof}

We use Vesztergombi's theorem in exactly the form of its main theorem: for any set of $m$ points in the plane, the second-largest distance occurs at most $\frac32 m$ times~\cite[Theorem]{Ve87}. There is no threshold on $m$ and no additive constant.

\begin{lemma}[(E)]\label{lem:E}
$\mult(\delta)\le e(S)+6r$.
\end{lemma}
\begin{proof}
Every $\delta$-pair not inside $S$ has an endpoint in $R$, and each point has at most $6$ points at distance $\delta$.
\end{proof}

\begin{lemma}\label{lem:lp}
Suppose $e(S)\le k s+C$ with $0\le k<3/2$. Then $\min\{\mult(\Dt),\mult(\delta)\}\le \frac{9}{7.5-k}\,n+C$.
\end{lemma}
\begin{proof}
Put $\lambda=(6-k)/(7.5-k)\in(0,1)$. By \cref{lem:VS,lem:E},
\[
\min\{\mult(\Dt),\mult(\delta)\}\le\lambda\cdot\tfrac32 s+(1-\lambda)(ks+C+6r)\le \tfrac{9}{7.5-k}(s+r)+C.
\]
For $k=1,\frac54,\frac43$ the factor $9/(7.5-k)$ is $\frac{18}{13},\frac{36}{25},\frac{54}{37}$.\chk{C1}
\end{proof}

\begin{lemma}\label{lem:lp2}
Suppose $e(S)\le s+2p+C$. Then $\min\{\mult(\Dt),\mult(\delta)\}\le\frac75 n+C$.
\end{lemma}
\begin{proof}
At most $p$ $\Dt$-pairs touch $Q$, and the rest lie inside $S\setminus Q$, so $\mult(\Dt)\le p+\frac32(s-p)$ by \cref{lem:VS}. Hence
$\min\le\frac45\bigl(\frac32 s-\frac12p\bigr)+\frac15\bigl(s+2p+C+6r\bigr)=\frac75 s+\frac65 r+\frac C5\le\frac75n+C$.\chk{C1}
\end{proof}

So in the nondegenerate regime it suffices to show $e(S)\le\frac43 s+O(1)$ in every $\tau$-regime except one, and to show that the remaining regime satisfies the hypothesis of \cref{lem:lp2}.

\subsection*{Small cases}
Suppose we are in the nondegenerate regime and $\Dt<10^4$. Then $\Delta<1.94\cdot10^4$, so $X$ lies in a disc of radius $\Delta$ about any of its points, and since the discs of radius $\frac12$ about points of $X$ are disjoint, $n\le(2\Delta+1)^2<N_0:=1.51\cdot10^9$.\chk{C11} Since $\mult(\delta)\le 3n$, Theorems~1 and~1$'$ hold for such $X$ once $C_0\ge 3N_0$. \emph{From now on, in the nondegenerate regime, $\Dt\ge 10^4$}; in particular $2/\Dt\le\beta/2$.\chk{C3}

\section{Bad edges are \texorpdfstring{$O(1)$}{O(1)} in the nondegenerate regime}\label{sec:bad}

\subsection{A lens lemma}

\begin{lemma}[lens]\label{lem:lens}
Let $\Lambda=\Dd(c_1,\rho)\cap\Dd(c_2,\rho)$ with $0<|c_1c_2|=d<2\rho$, and put $\cos\varphi_0:=d/(2\rho)$. Assume
\begin{equation}\label{eq:lenshyp}
\rho\ \ge\ \frac{2}{\cos(\varphi_0/2)-\cos\varphi_0}.
\end{equation}
Let $p$ lie on the arc of $\partial\Lambda$ belonging to $\Cc(c_1,\rho)$, let $q$ lie on the arc belonging to $\Cc(c_2,\rho)$, and suppose $|pq|\le 1$. Then there is a corner $c$ of $\Lambda$ such that the boundary arcs from $p$ to $c$ and from $c$ to $q$ each have length at most $\rho_0:=1/\sin(\varphi_0/2)$.
\end{lemma}
\begin{proof}
Take $c_1=(-d/2,0)$, $c_2=(d/2,0)$, $p=c_1+\rho(\cos\varphi,\sin\varphi)$ and $q=c_2+\rho(-\cos\psi,\sin\psi)$ with $|\varphi|,|\psi|\le\varphi_0$.
\emph{Step 1.} $p_x-q_x=\rho(\cos\varphi-\cos\varphi_0)+\rho(\cos\psi-\cos\varphi_0)$ with both terms $\ge0$; since $|p_x-q_x|\le1$, each term is at most $1$. By~\eqref{eq:lenshyp}, $\cos\varphi<\cos(\varphi_0/2)$, i.e.\ $|\varphi|>\varphi_0/2$; likewise $|\psi|>\varphi_0/2$.
\emph{Step 2.} By the mean value theorem, $\rho(\varphi_0-|\varphi|)=\rho(\cos\varphi-\cos\varphi_0)/\sin\xi\le 1/\sin\xi$ for some $\xi\in(\varphi_0/2,\varphi_0)\subseteq(0,\pi/2)$, so $\rho(\varphi_0-|\varphi|)\le\rho_0$; likewise for $\psi$.
\emph{Step 3.} If $\varphi,\psi$ had opposite signs, then $|p_y-q_y|\ge 2\rho\sin(\varphi_0/2)$. Now~\eqref{eq:lenshyp} gives $\rho\ge 2/(1-\cos\varphi_0)=1/\sin^2(\varphi_0/2)$, so $2\rho\sin(\varphi_0/2)\ge 2/\sin(\varphi_0/2)\ge2>1$, a contradiction.
\emph{Step 4.} So the corner $c=(0,\operatorname{sign}(\varphi)\rho\sin\varphi_0)$ works.
\end{proof}

\begin{lemma}[lens-corner arc length]\label{lem:lenscorner}
In the situation of \cref{lem:lens}, suppose moreover that $p,q\in\partial K'$ for a compact convex $K'\subseteq\Lambda$, and $|pq|=1$. Let $H$ be the closed half-plane bounded by the line $pq$ that contains $c$, and let $\gamma:=\partial K'\cap H$ be the boundary arc of $K'$ from $p$ to $q$ on $c$'s side. Then $\operatorname{length}(\gamma)\le 2\rho_0$.
\end{lemma}
\begin{proof}
$\partial\Lambda$ contains no segment, so the line $pq$ meets it only in $p$ and $q$, and $\partial\Lambda\cap H$ is exactly the arc $p\to c\to q$, of length at most $2\rho_0$. Since $K'\cap H\subseteq\Lambda\cap H$ are convex, monotonicity of perimeter gives $\operatorname{length}(\gamma)+1\le 2\rho_0+1$. (If $K'\cap H$ is the segment $[p,q]$, then $\gamma$ is that segment, of length $1\le2\rho_0$.)
\end{proof}

\emph{Constants.}\chk{C7} For \emph{KK lenses} ($\rho=\Dt$, $d\le\Delta\le1.94\Dt$) we have $\cos\varphi_0\le 0.97$, so $\rho_0^2=2/(1-\cos\varphi_0)\le 66.7<8.2^2$. The function $c\mapsto\cos(\varphi_0/2)-c=\sqrt{(1+c)/2}-c$ is decreasing on $[0,1)$, so~\eqref{eq:lenshyp} holds as soon as $\rho\ge 89.01$, in particular for $\rho=\Dt\ge10^4$. For \emph{DD lenses} ($\rho=\Delta$, $d\le\Delta$) we have $\cos\varphi_0\le\frac12$, $\rho_0\le 2$, and~\eqref{eq:lenshyp} needs only $\rho\ge 5.5$.

\subsection{The count}
We split bad edges $pq$ by type: \emph{KK} (both ends in $S\setminus D$), \emph{DD} (both in $S\cap D$) and \emph{cross} (one end in each).

If the perimeter of $K$ is below $1202$, then $|S\setminus D|<1202$ and at most $6\cdot1202$ edges touch $S\setminus D$; similarly for $P$ and $S\cap D$. We may therefore assume both perimeters are at least $1202$ when counting the bad edges touching the corresponding row; these small cases contribute a constant below $1.5\cdot10^4$ in total.

\subsubsection*{KK edges}
Let $p,q\in S\setminus D$ form a bad edge. If $w_p=w_q$, then $\angle(u_p,u_q)\le2\arcsin(1/(2\Dt))<\beta$,\chk{C3} so the edge would be good; hence $w_p\neq w_q$. Then $K\subseteq\Lambda:=\Dd(w_p,\Dt)\cap\Dd(w_q,\Dt)$, $p$ lies on the $\Cc(w_p,\Dt)$-arc and $q$ on the $\Cc(w_q,\Dt)$-arc of $\partial\Lambda$, and $d=|w_pw_q|\le\Delta\le1.94\Dt$. \Cref{lem:lens,lem:lenscorner} give an arc $\gamma_{pq}\subseteq\partial K$ from $p$ to $q$ of length at most $2\rho_0$. The union of the normal cones along the closed arc $\gamma_{pq}$ is connected and contains $u_p$ and $u_q$, so $\turn(\gamma_{pq})\ge\angle(u_p,u_q)>\beta$. If $y\in\gamma_{pq}$ then $p$ lies within arc distance $2\rho_0$ of $y$; there are at most $4\rho_0+1$ such $p$, each of degree at most $6$, so every point lies on at most $M=6(4\rho_0+1)$ of the arcs. By \cref{lem:multiplicity},
\[
\#\text{KK-bad}\le 12\pi(4\rho_0+1)/\beta<1.28\cdot10^5.\chk{C9}
\]

\subsubsection*{DD edges}
The same argument, on $\partial P$ with the lens $\Dd(x_p,\Delta)\cap\Dd(x_q,\Delta)$ and $\rho_0\le 2$, gives $\#\text{DD-bad}\le 12\pi\cdot 9/\beta<3.4\cdot10^4$. (If $x_p=x_q$ the angle is less than $2/\Delta<\beta$ and the edge is good.)

\subsubsection*{Cross edges}
Let $zv$ be a bad edge with $z\in S\cap D$ and $v\in S\setminus D$. Since $z\in X$, $|zw_v|\in[0,\Dt]\cup\{\Delta\}$.

\emph{(c1) $|zw_v|=\Delta$.} By \cref{lem:identities}(a), $(w_v-z)/\Delta$ is a second inward normal of $P$ at $z$. Its angle to $u_v$ is the angle at $w_v$ of the triangle $zvw_v$; that angle is opposite the side of length $1$, hence acute, and by the law of sines it is at most $\arcsin(1/\Delta)<2/\Dt\le\beta/2$.\chk{C3} As $\angle(u_z,u_v)>\beta$, the normal cone of $P$ at $z$ has angle greater than $\beta/2$. Normal cones at distinct vertices have disjoint interiors, so fewer than $4\pi/\beta$ points $z$ qualify, each of degree at most $6$: $\#(\mathrm{c1})<24\pi/\beta<7600$.

\emph{(c2) $|zw_v|\le\Dt$.} Let $g$ be the unit vector from $v$ towards $x_z$, and $\alpha:=\angle(u_v,g)$.
\begin{enumerate}[leftmargin=2em,label=(\roman*)]
\item \emph{Depth.} $v\notin D$ gives $|vx_z|\le\Dt$, while $|zx_z|=\Delta$ and $|zv|=1$ give $|vx_z|\ge\Delta-1$. So $v\in\Dd(x_z,\Dt)$ at depth $d_0:=\Dt-|vx_z|\in[0,1-t]$; in particular this case needs $t\le 1$. Also $K\subseteq\Dd(x_z,\Dt)$.
\item \emph{$\alpha$ is not small.} $\angle(u_z,g)$ is the angle at $x_z$ of the triangle $zvx_z$, which is at most $\arcsin(1/\Delta)<\beta/2$; hence $\alpha\ge\angle(u_v,u_z)-\angle(u_z,g)>\beta/2$.
\item \emph{$\alpha$ is not close to $\pi$.} Since $w_v,x_z\in X$, $|w_vx_z|\le\Delta\le1.94\Dt$. The law of cosines in the triangle $vw_vx_z$ (sides $\Dt$ and $\Dt-d_0$ at $v$, enclosing the angle $\alpha$) gives $\cos\alpha\ge(1+\varrho^2-1.94^2)/(2\varrho)$, where $\varrho:=(\Dt-d_0)/\Dt\in[1-10^{-4},1]$. The right side is increasing in $\varrho$, so $\cos\alpha\ge-0.8821$, i.e.\ $\alpha\le151.9^\circ$. By concavity of $\sin$ on $[0,\pi]$, $\sin\alpha\ge\min\{\sin(\beta/2),0.471\}=\sin(\beta/2)$.\chk{C8}
\item \emph{The walk.} Let $e_0\perp u_v$ be the unit vector with $e_0\cdot g=-\sin\alpha<0$. Follow $\partial K$ from $v$ in the direction whose forward one-sided tangent is closest to $e_0$, parametrised by arc length, $\gamma(0)=v$. Suppose the turning of $\gamma([0,\sigma])$ is at most $\beta/4$. This turning includes the whole normal cone $N_K(v)$. Then every normal along the arc is within $\beta/4$ of $u_v$, so every forward tangent $\gamma'$ is within $\beta/4$ of $e_0$ (at a corner $v$, the forward tangent is perpendicular to an extreme normal of $N_K(v)$, which is also within $\beta/4$ of $u_v$). Hence $\gamma'\cdot g\le-\sin\alpha+2\sin(\beta/8)\le-\sin(\beta/2)+2\sin(\beta/8)<-\beta/5$.\chk{C8} Integrating, $(\gamma(\sigma)-v)\cdot g\le-\sigma\beta/5$, and so
\[
|\gamma(\sigma)-x_z|^2\ \ge\ (\Dt-d_0)^2+\tfrac25(\Dt-d_0)\,\sigma\beta .
\]
For $\sigma=L:=6/\beta=600$ the right side exceeds $\Dt^2$,\chk{C8} which puts $\gamma(L)$ outside $\Dd(x_z,\Dt)\supseteq K$, a contradiction. Hence the closed arc $J_v$ of length $L$ from $v$ in the chosen direction has turning greater than $\beta/4$.
\item \emph{Count.} A point $y$ lies in $J_v$ only if $v$ is within arc distance $L$ of $y$, which allows at most $2L+1$ points $v$. By \cref{lem:multiplicity}, $\#\{v\}\le 2\pi(2L+1)/(\beta/4)$. Each such $v$ has degree at most $6$, so $\#(\mathrm{c2})\le48\pi(2L+1)/\beta<1.82\cdot10^7$.\chk{C9}
\end{enumerate}

\begin{proposition}\label{prop:bad}
In the nondegenerate regime the number $C_{\mathrm{bad}}$ of bad edges satisfies $C_{\mathrm{bad}}<1.83\cdot10^7$. Together with the constants $71$ and $12$ of \cref{sec:regions}, $C_{\mathrm{bad}}+71+12<10^8$.\chk{C9}
\end{proposition}

\section{The nondegenerate regime}\label{sec:regions}

\subsection{Region I: \texorpdfstring{$\tau\le\sqrt3/2-2\beta$}{tau <= sqrt3/2 - 2 beta}}\label{sec:regionI}
Let $p\in S\setminus D$ and let $q$ be a good neighbour of $p$; the \emph{elevation} of $q-p$ is $\arcsin((q-p)\cdot u_p)$. If $q$ is normal for $p$, then $b>0$ (by \cref{lem:identities}(a) if $q\in D$, by (b) if $q\notin D$), so $a<\beta$. Otherwise $q$ is an exception and $a=-\tau$.
\begin{itemize}[leftmargin=2em]
\item If $\tau\ge0$, every good neighbour on one side of $\R u_p$ has elevation in $[-\arcsin\tau,\arcsin\beta)$, a window of width less than $\arcsin(\sqrt3/2-2\beta)+\arcsin\beta\approx58.35^\circ<60^\circ$.\chk{C4}
\item If $\tau<0$, then $|\tau|\le1/(2\Dt)$ and all elevations lie in $(0,\arcsin\beta)\cup\{\arcsin|\tau|\}$; the window is again less than $60^\circ$.\chk{C4}
\end{itemize}
Let $p\in S\cap D$. Then $a>0$ always, and either $a<\beta$, or $p$ is the exception of $q$ and $a\in[\tau-\beta,\tau+\beta]$; the window is at most $\arcsin(\sqrt3/2-\beta)\approx58.87^\circ<60^\circ$.\chk{C4}

Since two neighbours of $p$ are at least $60^\circ$ apart, $p$ has at most one good neighbour on each side. Hence $e(S)\le s+C_{\mathrm{bad}}$, and \cref{lem:lp} with $k=1$ gives $\min\{\mult(\Dt),\mult(\delta)\}\le\frac{18}{13}n+C$.

\subsection{The regions with \texorpdfstring{$\tau>\sqrt3/2-2\beta$}{tau > sqrt3/2 - 2 beta}}
Here $\tau>2\beta$,\chk{C5,C6} so \cref{lem:dichotomy} applies. \emph{Region II} is $\tau\ge\sqrt3/2+3\beta$; \emph{Region III} is $\sqrt3/2-2\beta<\tau<\sqrt3/2+3\beta$.

\subsubsection*{Region II with $\tau>1$}
Then $t>1$, and by \cref{lem:identities}(c) there are no rungs. By (T1) the T-edges number at most $s$, so $e(S)\le s+C_{\mathrm{bad}}$, which gives $\frac{18}{13}n+C$.

From now on in Region II assume $\sqrt3/2+3\beta\le\tau\le1$. Then $t\le 1$ and $t\ge\tau$, so $t-\arcsin\beta\ge 0.886$; also $\sigma=\sqrt{1-\tau^2}\le\sigma_{\max}\approx0.4440$.\chk{C5}

\subsubsection{Mixed T-edges}
\begin{lemma}\label{lem:mixedT}
Assume $\sqrt3/2-2\beta<\tau\le1$. Let $v\in S\setminus D$ and $z\in S\cap D$ be joined by a T-edge. Then the closed arc of $\partial K$ of length $2$ starting at $v$ on $z$'s side has turning at least $t-\arcsin\beta$.
\end{lemma}
\begin{proof}
Suppose the turning is $\Theta<t-\arcsin\beta$; note that $\tau\le1$ gives $t\le1$, so $t-\arcsin\beta<1<\pi/3$. Let $e_0$ be the unit tangent direction at $v$ perpendicular to $u_v$ pointing to $z$'s side. Every forward tangent along the arc is within $\Theta$ of $e_0$, so $|\gamma(\sigma)-v|\ge\sigma\cos\Theta\ge\sigma/2$, and some $\sigma^*\le2$ has $|\gamma(\sigma^*)-v|=1$. The unit vector $\gamma(\sigma^*)-v$ lies within $\Theta$ of $e_0$, and the unit vector $z-v$ lies within $\arcsin\beta$ of $e_0$. Hence $|\gamma(\sigma^*)-z|\le\Theta+\arcsin\beta<t\le\dist(z,K)$ by \cref{lem:ballpoly}(c). But $\gamma(\sigma^*)\in K$, a contradiction.
\end{proof}
Each $v$ has at most two T-edges, and $v$ must lie within arc distance $2$ of a point $y$ of its arc, so at most $5$ such $v$ and $M\le10$. By \cref{lem:multiplicity} the number of mixed T-edges is at most $2\pi\cdot10/0.886<71$.\chk{C5} (\Cref{lem:mixedT} is stated for the larger range $\sqrt3/2-2\beta<\tau\le1$ because \cref{sec:loss1} uses it in Region~III. There $\tau<1$, so $t<1$ and $t\ge\tau>0.846$, hence $t-\arcsin\beta>0.836$, and the same count gives fewer than $2\pi\cdot10/0.836<76$ mixed T-edges; this is checked in \texttt{r4\_regionIII\_checks.py}.)

From here on, the \emph{K-row} is $S\setminus D$ and the \emph{D-row} is $S\cap D$, each with its own \emph{pure} T-edges; rungs join the two rows.

\subsubsection{Rungs}
For a rung $(v,z)$ write $z=v+o_{\sigma'}(u_v)$, where $o_{\pm}(u):=-\tau u\pm\sigma u^\perp$ (\cref{lem:identities}(c)). The map $u\mapsto o_\pm(u)$ is a rotation, so $|o_\pm(u)-o_\pm(u')|=|u-u'|\le\angle(u,u')$; and $|o_+(u)-o_-(u)|=2\sigma$. Hence $|o_{\pm}(u)-o_{\pm}(u')|\le2\sigma+\angle(u,u')$ for any choice of signs.

\begin{itemize}[leftmargin=2em]
\item[(R1)] \emph{Each vertex has at most one rung.} At $v$: the two candidate positions $v+o_\pm(u_v)$ are $2\sigma<1$ apart. At $z$: if $z$ had rungs to $v\neq v'$, both rungs are good, so $\angle(u_v,u_{v'})\le2\beta$ and $|v-v'|\le2\sigma+2\beta\approx0.908<1$,\chk{C5} a contradiction.
\end{itemize}

\emph{T-steps.} If $q$ is the right T-neighbour of $p$, then $q-p=\cos\eps\,u_p^\perp+\sin\eps\,u_p$ with $\eps\in(0,\arcsin\beta)$. So for T-steps $e_1$ at $v_0$ and $e_2$ at $z_0$, $\angle(e_1,e_2)\le\angle(u_{v_0},u_{z_0})+\arcsin\beta$.

\begin{lemma}[(R2), proximity]\label{lem:R2}
Let $(v_0,z_0)$ be a rung, $v_1$ the right pure T-neighbour of $v_0$ and $z_1$ the right pure T-neighbour of $z_0$; put $e_1=v_1-v_0$, $e_2=z_1-z_0$. If $v_1$ has a rung $(v_1,z')$, then $z'=z_1$.
\end{lemma}
\begin{proof}
$z'-z_1=(e_1-e_2)+(o_{\sigma_1}(u_{v_1})-o_{\sigma_0}(u_{v_0}))$. Since $\angle(u_{v_0},u_{z_0})\le\beta$ (good rung) and $\angle(u_{v_0},u_{v_1})\le\beta$ (good T-edge), $|z'-z_1|\le(\beta+\arcsin\beta)+\beta+2\sigma\approx0.918<1$.\chk{C5} Distinct points of $X$ are at least $1$ apart.
\end{proof}

\begin{lemma}[(R2-sym)]\label{lem:R2sym}
In the setting of \cref{lem:R2}, if $z_1$ has a rung $(v',z_1)$, then $v'=v_1$.
\end{lemma}
\begin{proof}
$v'=z_1-o_{\sigma'}(u_{v'})$ and $v_1=z_0-o_{\sigma_0}(u_{v_0})+e_1$, so $v'-v_1=(e_2-e_1)-(o_{\sigma'}(u_{v'})-o_{\sigma_0}(u_{v_0}))$. The angle between $u_{v'}$ and $u_{v_0}$ is at most $3\beta$, as the sum of three good-edge angles along $v'\to z_1\to z_0\to v_0$. Hence $|v'-v_1|\le(\beta+\arcsin\beta)+3\beta+2\sigma\le2\sigma+6\beta\approx0.948<1$.\chk{C5}
\end{proof}

\begin{itemize}[leftmargin=2em]
\item[(R3)] \emph{Parallel rungs and rhombi.} Rungs with the same vector $c$ satisfy $c\cdot u_v=-\tau$, so $u_v$ is one of the two unit solutions, which lie at $\pm\theta_0$ from $-c$, where $\theta_0:=\arccos\tau\in[0,\pi/2]$. Rungs form a matching by (R1), so their K-ends are distinct and by \cref{lem:nocommon} have distinct normals. Hence \emph{no three rungs are parallel}, and two parallel rungs have normals exactly $2\theta_0$ apart. If $v_0v_1$, $v_1z_1$, $z_1z_0$, $z_0v_0$ are all unit edges (a \emph{rhombus}), then $v_1$ and $z_0$ are the two points of $\Cc(v_0,1)\cap\Cc(z_1,1)$, so $v_0+z_1=v_1+z_0$ and the rungs $(v_0,z_0)$ and $(v_1,z_1)$ are parallel. Since the T-edge $v_0v_1$ is good, $2\theta_0=\angle(u_{v_0},u_{v_1})\le\beta$. So \emph{rhombi exist only if $\tau\ge\cos(\beta/2)$}, and $\cos(\beta/2)>\sqrt3/2+3\beta$.\chk{C5}
\end{itemize}

\subsubsection{Ownership and pairing}
For a rung $\mathsf r=(v,z)$, walk right from $v$ along pure K-row T-edges until reaching the first vertex after $v$ that is a rung vertex, or until the walk stops. Let $k\ge1$ be the number of vertices visited before stopping, $v$ included. Define $d$ in the same way in the D-row, starting from $z$. Then $\mathsf r$ \emph{owns} these $k+d$ vertices. Every other row vertex is \emph{unowned}: it has no rung vertex to its left on its T-path or T-cycle. Ownership is well defined, since each owned vertex belongs to the nearest rung vertex at or to its left, and that rung vertex lies on exactly one rung.

\emph{Charging.} Charge to $\mathsf r$ the rung itself and the right T-edge out of each vertex it owns. Charge every other pure T-edge to its unowned left endpoint. Every rung and pure T-edge is charged exactly once, and every vertex has at most one right T-edge. Hence:
\begin{itemize}[leftmargin=2em]
\item $\mathrm{edges}(\mathsf r)\le1+k+d$, with equality only if both walks end at a rung vertex (``both reach''); otherwise $\mathrm{edges}(\mathsf r)\le k+d$;
\item unowned vertices carry at most one edge each.
\end{itemize}

\emph{Both-reach gaps are $(1,1)$ or $(\ge2,\ge2)$.} If both walks reach and $k=1$, then $z_0$ has a right neighbour $z_1$ and $v_1$ is a rung vertex; by \cref{lem:R2} its rung is $(v_1,z_1)$, so $d=1$. If both reach and $d=1$, then $v_0$ has a right neighbour $v_1$ and $z_1$ is a rung vertex; by \cref{lem:R2sym} its rung is $(v_1,z_1)$, so $k=1$. A $(1,1)$ gap is a rhombus $v_0v_1z_1z_0$ (by (R1), $z_1\neq z_0$ and $v_1\neq v_0$). A cycle with a single rung vertex gives a both-reach gap with $k$ equal to the cycle length, which is at least $3$.

\emph{Ratios.} A rhombus gap has ratio edges/vertices $=3/2$; a both-reach gap with $k,d\ge2$ has ratio at most $5/4$; every other gap has ratio at most $1$.

If $\tau<\cos(\beta/2)$ there are no rhombi, so $e(S)\le\frac54s+C_{\mathrm{bad}}+71$ and \cref{lem:lp} gives $\frac{36}{25}n+C$.

If $\tau\ge\cos(\beta/2)$, pair each rhombus rung $\mathsf r_0=(v_0,z_0)$ with its successor $\mathsf r_1=(v_1,z_1)$. The rung $\mathsf r_1$ is not itself a rhombus rung: otherwise $\mathsf r_0,\mathsf r_1,\mathsf r_2$ would be three parallel rungs (here $v_2\neq v_0$ because T-cycles have length at least $3$). The pairing is injective, because $v_0$ is the unique left T-neighbour of $v_1$. If $\mathsf r_1$ owns $m$ vertices, the pair has ratio at most
\[
\max\Bigl\{\frac{4+m}{2+m}: m\ge4\Bigr\}\cup\Bigl\{\frac{3+m}{2+m}: m\ge2\Bigr\}=\frac43 .\chk{C5}
\]
Hence $e(S)\le\frac43s+C_{\mathrm{bad}}+71$, and \cref{lem:lp} gives $\frac{54}{37}n+C$. This is the only place where $54/37$ is attained.

\subsection{Region III: \texorpdfstring{$\sqrt3/2-2\beta<\tau<\sqrt3/2+3\beta$}{sqrt3/2 - 2 beta < tau < sqrt3/2 + 3 beta}}\label{sec:regionIII}
Here $\theta_0=\arccos\tau\in(26.36^\circ,32.22^\circ)$.

\begin{lemma}[corner lemma]\label{lem:corner}
Suppose $v\in S\setminus D$ has a rung partner $z$ and a second $\Dt$-partner $w'\neq w_v$, and put $u':=(w'-v)/\Dt$. Then $\angle(u_v,u')\ge\min\{2\theta_0,\,90^\circ-\theta_0\}>2\pi/7$.
\end{lemma}
\begin{proof}
Put $c=z-v$. If $|zw'|=\Delta$, then \cref{lem:identities}(c) applied with $w'$ gives $c\cdot u'=-\tau$, so $u'$ is the other unit solution and the angle is $2\theta_0$. If $|zw'|\le\Dt$, then $c\cdot u'>0$ while $\angle(c,u_v)=180^\circ-\theta_0$, so the angle is at least $90^\circ-\theta_0$. Numerically, $2\theta_0>2\pi/7$ because $\tau<\sqrt3/2+3\beta<\cos(\pi/7)$, and $90^\circ-\theta_0>2\pi/7$ because $\tau>\sqrt3/2-2\beta>\cos(3\pi/14)$.\chk{C6}
\end{proof}
Both $u_v$ and $u'$ are inward normals of $K$ at $v$, and $\angle(u_v,u')<\pi$ because $|w_vw'|\le\Delta<2\Dt$. So the normal cone $N_K(v)$ has angle greater than $2\pi/7$, and at most $6$ such $v$ exist. Every other K-vertex that has a rung has $\Dt$-degree exactly $1$, i.e.\ lies in $Q$. A K-vertex has at most $2$ rungs (the two positions $v+o_\pm(u_v)$), and (T1) bounds all T-edges, mixed ones included, by $s$. Hence
\[
e(S)\le s+2p+12+C_{\mathrm{bad}},
\]
and \cref{lem:lp2} gives $\frac75n+C$.

\begin{remark}
This regime genuinely needs the weighted \cref{lem:lp2}: a ``hexagonal strip'' around a single centre has average $\delta$-degree $3$ locally.
\end{remark}

\section{The degenerate regime \texorpdfstring{$\Delta>1.94\Dt$}{Delta > 1.94 Delta2}}\label{sec:degenerate}

\begin{proposition}\label{prop:degenerate}
If $\Delta>1.94\,\Dt$ and $n\ge4$, then $\mult(\Dt)\le n+1$.
\end{proposition}

Put $\kappa:=(\Dt/\Delta)^2<1/1.94^2$. Note that this section does not use $\delta$.

\emph{Step 1: no two disjoint diametral pairs.} Suppose $(x,x')$ and $(z,z')$ are disjoint diametral pairs with midpoints $m,m'$. The four cross distances satisfy $\sum|ab|^2=2\Delta^2+4|m-m'|^2\ge2\Delta^2$, and each is either at most $\Dt$ or equal to $\Delta$. If at most one equals $\Delta$, the sum is at most $\Delta^2(1+3\kappa)<2\Delta^2$,\chk{C10} a contradiction. So at least two cross distances equal $\Delta$.
\begin{itemize}[leftmargin=2em]
\item If two disjoint cross pairs equal $\Delta$, say $|xz|=|x'z'|=\Delta$, then $z$ and $x'$ are the two points of $\Cc(x,\Delta)\cap\Cc(z',\Delta)$, so $x,z,z',x'$ form a rhombus of side $\Delta$ whose diagonals satisfy $d_1^2+d_2^2=4\Delta^2$. One diagonal is at least $\sqrt2\Delta$, a contradiction.
\item If two cross pairs sharing a point equal $\Delta$, say $|xz|=|xz'|=\Delta$, then $xzz'$ is an equilateral triangle of side $\Delta$, and $x'\in\Cc(x,\Delta)$ with $|x'z|,|x'z'|\le\Delta$ lies on the $60^\circ$ arc of $\Cc(x,\Delta)$ between $z$ and $z'$, strictly inside it. One of $|x'z|,|x'z'|$ is at least $2\Delta\sin15^\circ$, and $(2\sin15^\circ)^2=2-\sqrt3>\kappa$,\chk{C10} so that distance lies in $(\Dt,\Delta)$, which is impossible.
\end{itemize}

\emph{Step 2: the diameter graph is a star.} A graph with no two disjoint edges is a star or a triangle. A triangle is impossible for $n\ge4$: a fourth point is not diametral, so it is within $\Dt$ of all three vertices, but $\Dt^2<\Delta^2/3$ is less than the squared circumradius.\chk{C10} So $D=\{y\}\cup A$ with $A$ the set of diametral partners of $y$; put $B:=X\setminus D\subseteq K$ and $W:=\{y\}\cup A$.

\emph{Step 3: where the $\Dt$-pairs are.}
\begin{itemize}[leftmargin=2em]
\item \emph{Within $B$: none.} For $a\in A$, $K\subseteq\Dd(y,\Dt)\cap\Dd(a,\Dt)$, a lens whose bounding box has squared diagonal less than $0.24\,\Dt^2$.\chk{C10} So $\operatorname{diam}K<\Dt$.
\item \emph{Within $A$: at most one.} $A\subseteq\Cc(y,\Delta)$, and its pairwise distances are at most $\Dt$ (a $\Delta$-pair inside $A$ would create a triangle in the diameter graph). A chord of $\Cc(y,\Delta)$ of length at most $\Dt<\Delta/1.94$ subtends an angle less than $30^\circ$ at $y$. So $A$ lies in an arc of angular width less than $60^\circ$, on which chord length is increasing in angular separation; a $\Dt$-pair in $A$ must realise the largest pairwise distance, hence be the two extreme points of $A$.
\item \emph{Between $y$ and $A$: none}, as those distances equal $\Delta$.
\end{itemize}
Hence $\mult(\Dt)\le1+I$, where $I:=\#\{(b,w):b\in B,\ w\in W,\ |bw|=\Dt\}$.

\emph{Step 4: $I\le|B|+|W|$.} We may assume $K$ has nonempty interior: an open chord between two points of $K$ lies in the interior of every disc $\Dd(w,\Dt)$, so otherwise $|B|\le|K\cap X|\le1$ and the bound is trivial. For $w\in X$ let $J_w:=\Cc(w,\Dt)\cap K$.
\begin{itemize}[leftmargin=2em]
\item \emph{Each $J_w$ is a single closed arc (possibly a point).} For $x\in X\setminus\{w\}$, $\Cc(w,\Dt)\cap\Dd(x,\Dt)$ is the closed arc $\{\cos(\theta-\theta_x)\ge|wx|/(2\Dt)\}$, of angular length at most $\pi$. For $p,q\in J_w$ we have $|pq|<\Dt<2\Dt$, so $p,q$ are not antipodal and the minor arc of $\Cc(w,\Dt)$ from $p$ to $q$ lies in every $\Dd(x,\Dt)$, hence in $K$. Since $\operatorname{diam}K<\Dt$, the angular spread of $J_w$ is less than $60^\circ$, so $J_w$ is a single closed arc.
\item \emph{No point is interior to two arcs.} Suppose $b$ is a relative-interior point of a nondegenerate arc $J_w$ and $b\in\Cc(w',\Dt)$ with $w'\in X\setminus\{w\}$. The function $f(\theta)=|w+\Dt e(\theta)-w'|^2-\Dt^2=|w-w'|^2+2\Dt\,e(\theta)\cdot(w-w')$ is $\le0$ on $J_w$ (as $K\subseteq\Dd(w',\Dt)$) and vanishes at $b$, so it has an interior local maximum with value $0$. But every local maximum of this shifted cosine is its global maximum $|w-w'|^2+2\Dt|w-w'|>0$, a contradiction.
\item \emph{Consequences.} Each nondegenerate $J_w$ is a subarc of $\partial K$, and distinct nondegenerate arcs have disjoint relative interiors. So a point $b\in B$ lies on at most two nondegenerate $W$-arcs (one on each side); any further incidences of $b$ come from degenerate arcs $J_w=\{b\}$.
\item \emph{Counting.} Let $I_{\mathrm{nd}}$ be the number of incidences from nondegenerate arcs. Then $I_{\mathrm{nd}}\le|B|+\#\{b:\text{two $W$-arcs meet at }b\}$. Mapping each such $b$ to the $W$-arc that starts at $b$ in the clockwise direction is injective, so $I_{\mathrm{nd}}\le|B|+\#\{w:J_w\text{ nondegenerate}\}$. Each degenerate arc contributes at most one incidence. Hence $I\le|B|+|W|$.
\end{itemize}
Therefore $\mult(\Dt)\le1+|B|+|A|+1=n+1$, proving \cref{prop:degenerate}.

\section{Assembly of the proof of Theorem 1\texorpdfstring{$'$}{'}}\label{sec:assembly}

\begin{table}[ht]
\centering
\renewcommand{\arraystretch}{1.2}
\begin{tabular}{@{}lll@{}}
\toprule
regime ($\delta=1$, $\beta=1/100$) & bound on $e(S)$ & bound on $\min\{\mult(\Dt),\mult(\delta)\}$\\
\midrule
degenerate: $\Delta>1.94\Dt$, $n\ge4$ & --- & $n+1$ (\cref{prop:degenerate})\\
nondegenerate, $\Dt<10^4$ & --- & $3n<3N_0$ (\cref{sec:reductions})\\
I: $\tau\le\sqrt3/2-2\beta$ (incl.\ $\tau\le0$) & $s+C$ & $\frac{18}{13}n+C$\\
III: $\sqrt3/2-2\beta<\tau<\sqrt3/2+3\beta$ & $s+2p+C$ & $\frac75n+C$\\
II: $\sqrt3/2+3\beta\le\tau<\cos(\beta/2)$ & $\frac54s+C$ & $\frac{36}{25}n+C$\\
II: $\cos(\beta/2)\le\tau\le1$ & $\frac43s+C$ & $\frac{54}{37}n+C$\\
II: $\tau>1$ ($\Leftrightarrow t>1$) & $s+C$ & $\frac{18}{13}n+C$\\
\bottomrule
\end{tabular}
\caption{The regimes in the proof of Theorem 1$'$; here $C<10^8$.}
\end{table}

The $\tau$-intervals partition $\R$, the split at $\Delta=1.94\,\Dt$ is exhaustive, and the cases $n\le3$ (or $X$ with a single distance) are trivial. Every constant in the last column is at most $54/37$,\chk{C1} and $n+1\le\frac{54}{37}n+1$. With $C<10^8$ from \cref{prop:bad}, Theorem~1$'$ holds with $C_0=3N_0+10^8$. \qed

\begin{remark}
The two least standard parts of the argument are the ownership/pairing bookkeeping in \cref{sec:regions} and the incidence count in Step~4 of \cref{sec:degenerate}. Both were re-derived independently, and both were also checked by computer in abstract models (\cref{app:verification}). The bookkeeping check is exhaustive up to $6+6$ row vertices, and both (R2-sym) and the ``no three parallel rungs'' fact turn out to be necessary for it. In this proof the bottleneck is the near-resonant regime $\tau\ge\cos(\beta/2)$, where rhombi occur. \Cref{sec:upgrade,sec:loss1,sec:loss2,sec:proof15} show that $54/37$ is not the limit of the method: the rungs counted above are paid for by points of $\Dt$-degree~$1$, and the term $6r$ of~(E) can be replaced by $5r+O(1)$.
\end{remark}

\section{From \texorpdfstring{$54/37$}{54/37} to \texorpdfstring{$15/11$}{15/11}}\label{sec:upgrade}

In \cref{sec:upgrade,sec:loss1,sec:loss2,sec:proof15} we keep the notation of \cref{sec:setup}; in particular $\delta=1$, $s=|S|$, $r=|R|$, and $Q$ is the set of points of $\Dt$-degree $1$, with $p=|Q|$. Except in \cref{sec:proof15}, we are in the nondegenerate regime $\Delta\le1.94\,\Dt$ and $\Dt\ge10^4$.

The proof of Theorem~1$'$ loses in two places.
\begin{itemize}[leftmargin=2em]
\item \emph{Loss 1: rungs are charged without their $\Dt$-cost.} Every $\delta$-edge that \cref{sec:regions} counts beyond $s$ is a rung: the rhombus rungs and the rungs of both-reach gaps in Region~II, and the second rungs at a K-vertex in Region~III. We show that, apart from $O(1)$ exceptions, the K-end of such a rung has $\Dt$-degree $1$, or else the rung forces a T-edge to be missing elsewhere. So $e(S)\le s+p+O(1)$ in every regime (\cref{sec:loss1}). The term $p$ is paid for on the $\Dt$ side, as in \cref{lem:lp2}.
\item \emph{Loss 2: the term $6r$ in~(E).} We replace~(E) by an exact identity and show that a point of $R$ contributes at most $5$ instead of $6$, apart from $O(1)$ points, and at most $4$ in Region~II (\cref{sec:loss2}).
\end{itemize}

\begin{lemma}[exact identity]\label{lem:identity}
For $y\in R$ let $\deg y$ be the number of points of $X$ at distance $1$ from $y$, and let $m_y$ be the number of those that lie in $S$. Then
\[
\mult(\delta)=e(S)+\sum_{y\in R}\frac{\deg y+m_y}{2}.
\]
\end{lemma}
\begin{proof}
There are $e(S)$ $\delta$-pairs inside $S$ and $\sum_{y\in R}m_y$ pairs with one end in $S$ and one in $R$. Each pair inside $R$ is counted twice in $\sum_{y\in R}(\deg y-m_y)$. Adding, $\mult(\delta)=e(S)+\sum_y m_y+\frac12\sum_y(\deg y-m_y)$.
\end{proof}
Since $\deg y\le6$ and $m_y\le\deg y$, \cref{lem:identity} contains~(E) (\cref{lem:E}) as the special case $\deg y+m_y\le12$.

\begin{lemma}[weighted linear programme]\label{lem:lp3}
Let $C,C'\ge0$.
\begin{enumerate}[label=(\alph*)]
\item If $e(S)\le s+p+C$ or $e(S)\le s+C$, and $\deg y+m_y\le10$ for all but at most $C'$ points $y\in R$, then $\min\{\mult(\Dt),\mult(\delta)\}\le\frac{15}{11}n+C+C'$.
\item If $e(S)\le s+p+C$ or $e(S)\le s+C$, and $\deg y+m_y\le8$ for all but at most $C'$ points $y\in R$, then $\min\{\mult(\Dt),\mult(\delta)\}\le\frac43n+C+2C'$.
\end{enumerate}
\end{lemma}
\begin{proof}
As in the proof of \cref{lem:lp2}, $\mult(\Dt)\le p+\frac32(s-p)=\frac32s-\frac12p$. In~(a), since $\deg y+m_y\le12$ always, \cref{lem:identity} gives $\mult(\delta)\le e(S)+5r+C'\le s+p+5r+C+C'$ (drop $p$ in the second case). Then
\[
\tfrac{8}{11}\bigl(\tfrac32s-\tfrac12p\bigr)+\tfrac{3}{11}\bigl(s+p+5r+C+C'\bigr)=\tfrac{15}{11}(s+r)-\tfrac1{11}p+\tfrac3{11}(C+C')\le\tfrac{15}{11}n+C+C'.
\]
Without the term $p$ in $e(S)$ the left side is $\frac{15}{11}(s+r)-\frac4{11}p+\frac{3}{11}(C+C')$, which is no larger. In~(b), $\mult(\delta)\le s+p+4r+C+2C'$, and the weights $\frac23,\frac13$ give $\frac43(s+r)+\frac13(C+2C')$, with coefficient $0$ on $p$. These certificates are checked in exact rational arithmetic in \texttt{r4\_lp\_checks.py}.
\end{proof}

\begin{remark}\label{rem:lpgeneral}
More generally, suppose $\deg y+m_y\le 2c$ for all but $O(1)$ points of $R$. If $e(S)\le s+O(1)$, the best linear combination gives $\frac{3c}{2c+1}n+O(1)$. If only $e(S)\le s+p+O(1)$ is known, the weight on $\mult(\Dt)$ must be at least $\frac23$ to cancel the term $p$, and the bound is $\max\{\frac43,\frac{3c}{2c+1}\}n+O(1)$. For $2c=12,10,8,6$ this gives $\frac{18}{13},\frac{15}{11},\frac43,\frac97$ in the first case and $\frac{18}{13},\frac{15}{11},\frac43,\frac43$ in the second (\texttt{r4\_lp\_checks.py}). In particular, with only~(E) and $e(S)\le s+p+C$, the weights $\frac{10}{13},\frac{3}{13}$ give $\frac{18}{13}s-\frac2{13}p+\frac{18}{13}r\le\frac{18}{13}n+C$; this is the intermediate bound obtained from \cref{sec:loss1} alone.
\end{remark}

\section{Loss 1: \texorpdfstring{$e(S)\le s+|Q|+O(1)$}{e(S) <= s + |Q| + O(1)}}\label{sec:loss1}

The aim of this section is the following bound.

\begin{proposition}\label{prop:eS}
In the nondegenerate regime with $\Dt\ge10^4$, in every $\tau$-regime,
\[
e(S)\le s+p+C_1,\qquad C_1\le 2C_{\mathrm{bad}}+100<3.7\cdot10^7 .
\]
In Region~I and for $\tau>1$ even $e(S)\le s+C_{\mathrm{bad}}$.
\end{proposition}

Region~I and the case $\tau>1$ were done in \cref{sec:regions}. Region~II ($\sqrt3/2+3\beta\le\tau\le1$) is \cref{prop:regionII2} and Region~III is \cref{prop:regionIII2}; together they prove \cref{prop:eS}, the largest constant being $2C_{\mathrm{bad}}+81$ in \cref{prop:tauexact}. Throughout this section $\tau>2\beta$, so \cref{lem:dichotomy} applies: every good edge is a T-edge or a rung. We write $e_T$ for the number of T-edges (pure and mixed) and $\varrho$ for the number of rungs. By (T1) every vertex has at most one T-neighbour on each side, so $e_T\le s$. We use $\theta_0:=\arccos\tau$, $\sigma=\sqrt{1-\tau^2}=\sin\theta_0$ and $o_\pm(u)=-\tau u\pm\sigma u^\perp$ as in \cref{sec:regions}.

\subsection{The corner lemma in general form}
The following lemma extends \cref{lem:corner} to all $\tau\in[0,1]$ and records which of its two cases occurs.

\begin{lemma}[corner lemma, general form]\label{lem:corner2}
Assume $0\le\tau\le1$. Let $v\in S\setminus D$ and $z\in X$ with $|vz|=1$ and $|zw_v|=\Delta$ (for instance, $z$ a rung partner of $v$), and put $c:=z-v$, so that $c\cdot u_v=-\tau$ by \cref{lem:identities}(c). Let $w'\neq w_v$ be a second $\Dt$-partner of $v$ and $u':=(w'-v)/\Dt$. Then exactly one of the following holds.
\begin{enumerate}[label=(\roman*)]
\item $|zw'|\le\Dt$. Then $c\cdot u'\ge1/(2\Dt)>0$, and $\angle(u_v,u')>\pi/2-\theta_0$.
\item $|zw'|=\Delta$. Then $u'=u^*$, the reflection of $u_v$ in the line $\R c$, and $\angle(u_v,u')=2\theta_0$.
\end{enumerate}
In both cases $u_v$ and $u'$ are inward normals of $K$ at $v$ with $\angle(u_v,u')<\pi$, so the normal cone of $K$ at $v$ has angle at least $\angle(u_v,u')$.
\end{lemma}
\begin{proof}
The cases are exhaustive because $|zw'|\in[0,\Dt]\cup\{\Delta\}$. \Cref{lem:identities} holds with any $\Dt$-partner of $v$ in place of $w_v$, since its proof uses only $|vw_v|=\Dt$.
(i) By \cref{lem:identities}(b) with $w'$ in place of $w_v$, $z$ is normal for $v$ with respect to $w'$, so $c\cdot u'\ge|c|^2/(2\Dt)=1/(2\Dt)$. Hence $\angle(c,u')<\pi/2$. As $|c|=1$ and $c\cdot u_v=-\tau=\cos(\pi-\theta_0)$, $\angle(c,u_v)=\pi-\theta_0$, and so $\angle(u_v,u')\ge\angle(c,u_v)-\angle(c,u')>\pi/2-\theta_0$.
(ii) By \cref{lem:identities}(c) with $w'$ in place of $w_v$, $c\cdot u'=-\tau$. The unit vectors $u$ with $c\cdot u=-\tau$ are $u_v$ and its reflection $u^*$ in $\R c$; both make the angle $\theta_0$ with $-c$, so $\angle(u_v,u^*)=2\theta_0$. Since $w'=v+\Dt u'\neq w_v=v+\Dt u_v$, $u'\ne u_v$, so $u'=u^*$ (and this case cannot occur when $\theta_0=0$).
Finally, $K\subseteq\Dd(w',\Dt)$ and $v\in\Cc(w',\Dt)$, so $u'$ is an inward normal of $K$ at $v$, like $u_v$; and $|w_vw'|\le\Delta<2\Dt$ gives $\angle(u_v,u')<\pi$.
\end{proof}

As in \cref{sec:regionIII}, normal cones of $K$ at distinct boundary points have disjoint interiors, so for any $\phi>0$ fewer than $2\pi/\phi$ points of $\partial K$ have a normal cone of angle greater than $\phi$. Let $v\in S\setminus D$ have a rung partner. We say $v$ is \emph{of type~(i)} if some $\Dt$-partner $w'\ne w_v$ of $v$ falls under case~(i) of \cref{lem:corner2} (for some rung partner $z$ of $v$). If $v$ has $\Dt$-degree at least $2$ and is not of type~(i), then all its partners other than $w_v$ lie in the single direction $u^*$ from $v$, so $v$ has $\Dt$-degree exactly $2$; we call such $v$ \emph{double}.

\subsection{Region II}

\begin{proposition}\label{prop:regionII2}
In Region~II, $\sqrt3/2+3\beta\le\tau\le1$, we have $e(S)\le s+p+C_{\mathrm{bad}}+5$.
\end{proposition}

Throughout this subsection $\sqrt3/2+3\beta\le\tau\le1$. Then $t\le1$ and $t\ge\tau$ (\cref{sec:setup}), and $\theta_0\le\arccos(\sqrt3/2+3\beta)<26.4^\circ$. By (R1) of \cref{sec:regions}, each vertex has at most one rung, so the rungs form a matching. Since $\pi/2-\theta_0>\pi/3$, at most $5$ K-vertices are of type~(i). Every other rung K-end either lies in $Q$ or is double.

\begin{lemma}[far partners]\label{lem:farpartner}
Let $z,w\in X$ with $|zw|=\Delta$, put $\hat w:=(w-z)/\Delta$, and let $y\in X$ with $|zy|=1$ and $\eta:=(y-z)\cdot\hat w$. Then $|yw|^2=\Delta^2+1-2\Delta\eta$. If
\[
\eta<\frac{\Delta^2-\Dt^2+1}{2\Delta}=t+\frac{1-t^2}{2\Delta},
\]
then $|yw|=\Delta$. When $t\le1$ the threshold is at least $t$.
\end{lemma}
\begin{proof}
The identity is the law of cosines. If $\eta$ is below the threshold, then $|yw|^2>\Dt^2$, and $X$ has no distance in $(\Dt,\Delta)$. The equality of the two forms of the threshold follows from $\Dt=\Delta-t$.
\end{proof}

\begin{lemma}[double rungs]\label{lem:doublerung}
Let $v$ be double with rung partner $z$. Then $z$ has no T-neighbour at all, pure or mixed.
\end{lemma}
\begin{proof}
Let $w:=w_v$, and let $w^*$ be the partner of $v$ in direction $u^*$. By the definition of a rung and by case~(ii) of \cref{lem:corner2}, $|zw|=|zw^*|=\Delta$. Put $\hat w:=(w-z)/\Delta$ and $\hat w^*:=(w^*-z)/\Delta$.
The angle between $\hat w$ and $u_v=(w-v)/\Dt$ is the angle at $w$ of the triangle $vzw$. That angle is opposite the side $vz$ of length $1$, the shortest side, so it is acute, and by the law of sines it is at most $\arcsin(1/\Delta)$. Likewise $\angle(\hat w^*,u^*)\le\arcsin(1/\Delta)$. Also $\angle(u^*,u_v)=2\theta_0$, and $\angle(u_v,u_z)\le\beta$ because the rung is good.
Suppose $z_1$ is a T-neighbour of $z$ and put $e:=z_1-z$. Then $e\cdot u_z\in(0,\beta)$, so $\angle(e,u_z)\ge\pi/2-\arcsin\beta$. With $\eps:=\beta+\arcsin\beta+\arcsin(1/\Delta)$ we get $\angle(e,\hat w)\ge\pi/2-\eps$ and $\angle(e,\hat w^*)\ge\pi/2-\eps-2\theta_0\ge0$, hence
\[
\eta:=e\cdot\hat w\le\sin\eps,\qquad \eta^*:=e\cdot\hat w^*\le\sin(2\theta_0+\eps).
\]
Now $\eps<1.2^\circ$ and $3\theta_0+\eps<3\cdot26.4^\circ+1.2^\circ<90^\circ$, so $\sin(2\theta_0+\eps)<\sin(\pi/2-\theta_0)=\tau\le t$, and also $\sin\eps<t$. By \cref{lem:farpartner}, $|z_1w|=|z_1w^*|=\Delta$. Since $u^*\ne u_v$, $w^*\ne w$. So $z$ and $z_1$ are two distinct points of $\Cc(w,\Delta)\cap\Cc(w^*,\Delta)$. These two points are mirror images in the line $ww^*$, at distance $2\sqrt{\Delta^2-|ww^*|^2/4}\ge\sqrt3\,\Delta>1$, because $|ww^*|\le\Delta$. This contradicts $|zz_1|=1$.
\end{proof}
The script \texttt{r4\_regionII\_checks.py} checks in exact rational arithmetic an alternative sufficient condition for the two inequalities $\eta,\eta^*<t$ over all of Region~II, and reports an adversarial floating-point search against \cref{lem:doublerung}.

\begin{proof}[Proof of \cref{prop:regionII2}]
Call a rung \emph{double} if its K-end is double. By \cref{lem:doublerung} the D-end of every double rung has no T-edge, and these D-ends are distinct because the rungs form a matching. Every vertex has at most two T-edges, so, counting T-edges by their endpoints, $e_T\le s-\#\{\text{double rungs}\}$. Every good edge is a T-edge or a rung, so
\[
e(S)\le e_T+\varrho+C_{\mathrm{bad}}\le s+\#\{\text{rungs that are not double}\}+C_{\mathrm{bad}}.
\]
The K-end of a rung that is not double is of type~(i) (at most $5$ of them) or lies in $Q$. Since each K-vertex has at most one rung, there are at most $p+5$ such rungs.
\end{proof}

For \cref{prop:regionII2} we do not need the mixed T-edge lemma, (R2), (R2-sym), the rhombi, or the ownership/pairing argument of \cref{sec:regions}.

\subsection{Region III}

\begin{proposition}\label{prop:regionIII2}
In Region~III, $\sqrt3/2-2\beta<\tau<\sqrt3/2+3\beta$, we have $e(S)\le s+p+2C_{\mathrm{bad}}+81$.
\end{proposition}

Throughout this subsection $\tau$ is in Region~III. Then $\tau<1$, so $t<1$ and $t\ge\tau$. By \cref{lem:mixedT} and the remark after it, there are fewer than $76$ mixed T-edges. A K-vertex $v$ has at most two rungs, at the positions $v+o_\pm(u_v)$. Let $q_1$ and $q_2$ be the numbers of K-vertices with exactly one and exactly two rungs. Then $\varrho\le q_1+2q_2$. By \cref{lem:corner} (equivalently \cref{lem:corner2}, since $\min\{2\theta_0,\pi/2-\theta_0\}>2\pi/7$ in Region~III), at most $6$ K-vertices with a rung have $\Dt$-degree at least $2$, so $q_1+q_2\le p+6$. Since $e(S)\le e_T+\varrho+C_{\mathrm{bad}}$, it is enough to show
\begin{equation}\label{eq:eTgoal}
e_T\le s-q_2+C_2,
\end{equation}
because then $e(S)\le s+q_1+q_2+C_2+C_{\mathrm{bad}}\le s+p+C_2+C_{\mathrm{bad}}+6$. We show~\eqref{eq:eTgoal} with $C_2=0$ for $\tau\ne\sqrt3/2$ and with $C_2=C_{\mathrm{bad}}+75$ for $\tau=\sqrt3/2$.

\subsubsection*{(a) $\sqrt3/2<\tau<\sqrt3/2+3\beta$}
Here $2\sigma<1$. The two possible rung positions $v+o_\pm(u_v)$ are $2\sigma<1$ apart, so they cannot both be points of $X$, and $q_2=0$. So~\eqref{eq:eTgoal} is $e_T\le s$, which holds by (T1).

\subsubsection*{(b) $\sqrt3/2-2\beta<\tau<\sqrt3/2$}
Here $0<2\sigma-1<0.067$.

\begin{lemma}[slot killing]\label{lem:slot}
Let $v$ have two rungs $z=v+o_+(u_v)$ and $z'=v+o_-(u_v)$. Then $z'$ has no T-neighbour (pure or mixed) on the side of $z$, and $z$ has no T-neighbour on the side of $z'$. Here ``the side of $z$'' at $z'$ is the side containing $z$, i.e.\ the sign of $(z-z')\cdot u_{z'}^\perp$.
\end{lemma}
\begin{proof}
We have $z-z'=2\sigma u_v^\perp$ and $\angle(u_v,u_{z'})\le\beta$ (the rung is good), so $\angle(u_v^\perp,u_{z'}^\perp)\le\beta$; in particular $z$ is on the right of $z'$. A right T-neighbour $y$ of $z'$ satisfies $y-z'=\cos\eps\,u_{z'}^\perp+\sin\eps\,u_{z'}$ with $\eps\in(0,\arcsin\beta)$. Hence $\angle(y-z',z-z')\le\gamma:=\beta+\arcsin\beta<0.0201$ and
\[
|y-z|^2=1+4\sigma^2-4\sigma\cos\angle(y-z',z-z')\le(2\sigma-1)^2+4\sigma(1-\cos\gamma)<0.005<1.
\]
Also $y\ne z$, since $|zz'|=2\sigma\ne1$. This contradicts $\delta=1$. The statement for $z$ is symmetric.
\end{proof}

Call a pair (vertex, side) a \emph{slot}; each vertex has two slots, and $e_T$ is half the number of slots used by T-edges. By \cref{lem:slot}, each K-vertex $v$ with two rungs \emph{kills} two slots: the slot of $z'$ on the side of $z$ and the slot of $z$ on the side of $z'$. These $2q_2$ killed slots are distinct. Indeed, if $x=v+o_-(u_v)$ then $v$ kills the right slot of $x$, and if $x=v+o_+(u_v)$ then $v$ kills the left slot of $x$ (right and left with respect to $u_x^\perp$, which is within $\beta$ of $u_v^\perp$). So if $v\ne\bar v$ both have two rungs and kill the same slot of a point $z'$, then $z'$ plays the same role for both; by symmetry say $z'=v+o_-(u_v)=\bar v+o_-(u_{\bar v})$. Let $z$ and $\bar z$ be the other rung partners of $v$ and $\bar v$. Then $\bar z=z'+2\sigma u_{\bar v}^\perp$ and $z=z'+2\sigma u_v^\perp$, and $\angle(u_v,u_{\bar v})\le2\beta$ (both rungs at $z'$ are good), so $|\bar z-z|\le2\sigma\cdot2\beta<1$ and $\bar z=z$. Then $v,\bar v$ are the two points of $\Cc(z,1)\cap\Cc(z',1)$, mirror images in the line $zz'$. From $v-z'=\tau u_v+\sigma u_v^\perp$ we get $\bar v-z'=-\tau u_v+\sigma u_v^\perp$ and so $(\bar v-z')\cdot u_{z'}\le-\tau\cos\beta+\sigma\sin\beta<0$, contradicting \cref{lem:identities}(a) at $z'\in S\cap D$. Hence at most $2s-2q_2$ slots are used, and $e_T\le s-q_2$, which is~\eqref{eq:eTgoal}. The constants in (a) and (b) are checked in exact rational arithmetic in \texttt{r4\_regionIII\_checks.py}; a floating-point search found $\max|y-z|=0.0691$, against the proved bound $\sqrt{0.005}<0.0708$.

\subsubsection*{(c) $\tau=\sqrt3/2$}
Now $\sigma=\frac12$ and $\theta_0=30^\circ$, and slot killing fails because $|zz'|=1$. We also note that $t<1$ gives $t<\sqrt3/2+1/(2\Dt)<0.8661$.

\begin{proposition}\label{prop:tauexact}
If $\tau=\sqrt3/2$, then $e_T\le s-q_2+C_{\mathrm{bad}}+75$.
\end{proposition}

\emph{Triangles.} Let $v$ have two rungs $z'=v+o_-(u_v)$ and $z=v+o_+(u_v)$. Then $z-z'=u_v^\perp$, so $vz'z$ is an equilateral triangle of side $1$, and $z'z$ is a $\delta$-edge between two points of $S\cap D$; we call it the \emph{base} and $v$ the \emph{apex}. Since $(z-z')\cdot u_{z'}^\perp=u_v^\perp\cdot u_{z'}^\perp\ge\cos\beta>0$, $z$ is on the right of $z'$ (if the base is a T-edge, $z$ is the right T-neighbour of $z'$).

\emph{A base determines its apex.} The other point of $\Cc(z,1)\cap\Cc(z',1)$ is $\bar v=z'-\tau u_v+\frac12u_v^\perp$, and $(\bar v-z')\cdot u_{z'}\le-\tau\cos\beta+\frac12\sin\beta<0$, which \cref{lem:identities}(a) forbids for a point of $X$. So distinct apexes have distinct bases, and at most $C_{\mathrm{bad}}$ triangles have a bad base. A good base joins two points of $S\cap D$, so by \cref{lem:dichotomy} it is a pure D-row T-edge.

\emph{Counting T-edges by left ends.} By (T1) and (T2), each pure T-edge has exactly one \emph{left end} (the endpoint of which the other endpoint is the right T-neighbour), and each vertex is the left end of at most one T-edge. Let $\mathrm{br}_K$ (resp.\ $\mathrm{br}_D$) be the number of K-row (resp.\ D-row) vertices that have no \emph{pure} right T-neighbour, i.e.\ no right T-neighbour in the same row. Then the numbers of pure T-edges in the two rows are $|S\setminus D|-\mathrm{br}_K$ and $|S\cap D|-\mathrm{br}_D$, and with fewer than $76$ mixed T-edges,
\begin{equation}\label{eq:breaks}
e_T\le s-\mathrm{br}_K-\mathrm{br}_D+75 .
\end{equation}
So it suffices to show $\mathrm{br}_K+\mathrm{br}_D\ge q_2-C_{\mathrm{bad}}$. We do this by finding one break per good-base triangle.

\emph{Orientation.} For $p\in S$ the body ($K$ or $P$) lies on the side of $u_p$, and $u_p^\perp$ is $u_p$ rotated by $-90^\circ$. So a step from $p$ in direction close to $u_p^\perp$ runs along the boundary counterclockwise. For points on a circle $\Cc(w,\cdot)$ whose direction field is within $O(\beta)$ of the direction towards $w$ (this is the case below for both rows), ``to the right'' means ``counterclockwise about $w$'', in both rows.

\begin{lemma}[one circle per component]\label{lem:onecircle}
Let $Z$ be the connected component of the graph of pure D-row T-edges that contains a good base $z'z$ of a triangle with apex $v$, and let $w:=w_v$. Then every vertex of $Z$ lies on $\Cc(w,\Delta)$; consecutive vertices of $Z$ are at angle exactly $\alpha_1:=2\arcsin(1/(2\Delta))$ about $w$, each right step turning counterclockwise about $w$; and $Z$ is a path spanning an angle at most $\pi/3$ about $w$.
\end{lemma}
\begin{proof}
By the definition of a rung, $|zw|=|z'w|=\Delta$.
\emph{Interior vertices.} Let $z_k\in Z$ have a left T-neighbour $z_{k-1}$ and a right T-neighbour $z_{k+1}$ in $Z$, and put $u:=u_{z_k}$. Then $z_{k+1}-z_k=\cos\eps_1u^\perp+\sin\eps_1u$, where $\sin\eps_1\in(0,\beta)$ is the $a$-value of the T-edge $z_kz_{k+1}$, and $z_{k-1}-z_k=-\cos\eps_2u^\perp+\sin\eps_2u$, where $\sin\eps_2\in(0,\beta)$ is the $b$-value of the T-edge $z_{k-1}z_k$. Let $\nu=\cos\phi\,u+\sin\phi\,u^\perp$ ($\phi\in(-\pi,\pi]$) be an inward normal of $P$ at $z_k$. Then $\nu\cdot(z_{k+1}-z_k)=\sin(\eps_1+\phi)\ge0$ and $\nu\cdot(z_{k-1}-z_k)=\sin(\eps_2-\phi)\ge0$, so $\phi\in[-\eps_1,\eps_2]$.
\emph{Induction step.} Suppose moreover $|z_kw|=\Delta$. Every $x\in X$ has $|xw|\le\Delta$, so, as in \cref{lem:identities}(a), $\hat w_k:=(w-z_k)/\Delta$ is an inward normal of $P$ at $z_k$. Hence $\eta:=(z_{k+1}-z_k)\cdot\hat w_k=\sin(\eps_1+\phi)\in[0,\sin(2\arcsin\beta)]$, and $\sin(2\arcsin\beta)<0.0201<t$. By \cref{lem:farpartner}, $|z_{k+1}w|=\Delta$. In the same way $(z_{k-1}-z_k)\cdot\hat w_k=\sin(\eps_2-\phi)\in[0,0.0201]$ gives $|z_{k-1}w|=\Delta$.
\emph{Induction.} Walk right from $z$: $z$ has the left neighbour $z'$, so if $z$ has a right neighbour it is interior and that neighbour lies on $\Cc(w,\Delta)$; the new vertex again has a left neighbour, and so on. Walk left from $z'$ in the same way. Every vertex of $Z$ is reached by one of the two walks, with no wrap-around needed.
\emph{Angles.} Consecutive vertices of $Z$ are points of $\Cc(w,\Delta)$ at distance $1$, so at angle $\alpha_1$ about $w$. Let $\hat w_k^\perp$ be $\hat w_k$ rotated by $-90^\circ$; a step $z_k\to z_{k+1}$ between points of $\Cc(w,\Delta)$ is counterclockwise about $w$ if and only if $(z_{k+1}-z_k)\cdot\hat w_k^\perp>0$, and this sense does not depend on the endpoint at which it is evaluated. If $z_k$ is interior, $(z_{k+1}-z_k)\cdot\hat w_k^\perp=\cos(\eps_1+\phi)>0$. If $z_{k+1}$ is interior, then in the same way $(z_k-z_{k+1})\cdot\hat w_{k+1}^\perp=-\cos(\eps_2-\phi)<0$ with the quantities of $z_{k+1}$. Every edge of $Z$ other than the base has an interior endpoint, and for the base, $z-z'=u_v^\perp$, where $u_v$ is the direction from the midpoint $\frac12(z+z')=v-\tau u_v$ of the base to $w=v+\Dt u_v$. So each right step turns counterclockwise about $w$. Two points of $X$ are at most $\Delta$ apart, and a chord $2\Delta\sin(\vartheta/2)\le\Delta$ has $\vartheta\le\pi/3$, so $Z$ lies in an arc of angle at most $\pi/3$ about $w$. Monotone steps of $\alpha_1$ inside such an arc cannot close up, so $Z$ is a path.
\end{proof}

\begin{lemma}[one centre per component]\label{lem:onecentre}
If two triangles with apexes $v_A,v_B$ have good bases in the same component $Z$, then $w_{v_A}=w_{v_B}$.
\end{lemma}
\begin{proof}
By \cref{lem:onecircle} for $v_A$, the base $z_B'z_B$ of the second triangle lies on $\Cc(w_{v_A},\Delta)$; by the definition of a rung it also lies on $\Cc(w_{v_B},\Delta)$. Suppose $w_{v_A}\ne w_{v_B}$; then $w_{v_A}$ is the mirror image of $w_{v_B}$ in the line $z_B'z_B$. Put $a:=(w_{v_B}-z_B)/\Delta$ and let $e$ be the unit vector along the base. Since $z_B'z_B$ is a chord of length $1$ of $\Cc(w_{v_B},\Delta)$, $a\cdot e=\pm1/(2\Delta)$ exactly. The mirror image of $a$ is $a^*=2(a\cdot e)e-a$, and $a\cdot a^*=2(a\cdot e)^2-1<0$, so $(w_{v_A}-z_B)\cdot a=\Delta\,a\cdot a^*<0$. But $w_{v_B}$ is a diametral partner of $z_B$ and $w_{v_A}\in X$, so $(w_{v_A}-z_B)\cdot a\ge|w_{v_A}-z_B|^2/(2\Delta)\ge0$ by the argument of \cref{lem:identities}(a). Contradiction.
\end{proof}

Fix a component $Z$ containing at least one good base, and let $w$ be the common centre. The apex $v_A$ of a triangle with base in $Z$ lies on $\Cc(w,\Dt)$, and since $v_A=w-\Dt u_{v_A}$, the midpoint of its base is $v_A-\tau u_{v_A}=w-(\Dt+\tau)u_{v_A}$. So $v_A$ and the midpoint of its base lie on the same ray from $w$ and have the same angle about $w$. Order the triangles with base in $Z$ by the position of the base along $Z$. Let $A,B$ be consecutive, $B$ to the right, with bases $(z_i,z_{i+1})$ and $(z_{i+k},z_{i+k+1})$ in the numbering of \cref{lem:onecircle}, $k\ge1$. Then the apexes are at angle $k\alpha_1$ about $w$, with
\begin{equation}\label{eq:kalpha}
\alpha_2:=2\arcsin\frac1{2\Dt}\ \le\ k\alpha_1\ \le\ \frac\pi3-\alpha_1<\frac\pi3,
\end{equation}
the lower bound because $|v_Av_B|\ge1$, the upper bound because $Z$ spans at most $\pi/3$. Let $\gamma_{AB}$ be the shorter arc of $\Cc(w,\Dt)$ from $v_A$ to $v_B$.

\begin{lemma}[the arc lies on $\partial K$]\label{lem:arcK}
$\gamma_{AB}\subseteq\partial K$.
\end{lemma}
\begin{proof}
For $x\in X\setminus\{w\}$, $\Dd(x,\Dt)\cap\Cc(w,\Dt)$ is the arc $\{\vartheta:\cos(\vartheta-\vartheta_x)\ge|wx|/(2\Dt)\}$, of angle at most $\pi$ ($\vartheta_x$ the direction of $x$ from $w$); for $x=w$ it is the whole circle. It contains $v_A,v_B\in K$, which are less than $\pi$ apart, so it contains $\gamma_{AB}$. Hence $\gamma_{AB}\subseteq K$, and as $K\subseteq\Dd(w,\Dt)$, $\gamma_{AB}\subseteq\partial K$.
\end{proof}

\begin{lemma}[rigidity]\label{lem:rigid}
Starting from $v_0:=v_A$, follow pure right T-edges in the K-row: $v_0,v_1,\dots$. As long as $v_i\in\gamma_{AB}$ and $v_i\ne v_B$, either $v_i$ has no pure right T-neighbour, or $v_{i+1}$ is the point of $\gamma_{AB}$ at angle $\alpha_2$ counterclockwise from $v_i$ about $w$.
\end{lemma}
\begin{proof}
\emph{Normal fact.} Put $\hat w_i:=(w-v_i)/\Dt$. Then $u_{v_i}=\hat w_i$. For $i=0$ this holds because $w=w_{v_A}$. If $v_i\ne v_A$, then $v_i$ is in the relative interior of $\gamma_{AB}\subseteq\partial K$, and an inward normal $\nu$ of $K$ at $v_i$ satisfies $(y-v_i)\cdot\nu\ge0$ for all $y$ of the arc on both sides of $v_i$; this forces $\nu$ to be the direction $\hat w_i$ of the circle's normal. Since $u_{v_i}$ is an inward normal of $K$ at $v_i$, $u_{v_i}=\hat w_i$.
\emph{Step.} Let $v^*$ be the point of $\Cc(v_i,1)\cap\Cc(w,\Dt)$ counterclockwise from $v_i$, at angle $\alpha_2$ about $w$. It lies on $\gamma_{AB}$: otherwise $v_B$ would be at angle in $(0,\alpha_2)$ from $v_i$, so $|v_iv_B|<1$. Suppose $v_{i+1}$ is the pure right T-neighbour of $v_i$. In coordinates with first axis $u_{v_i}^\perp$ and second axis $u_{v_i}=\hat w_i$, write $v_{i+1}-v_i=(\cos\psi,\sin\psi)$ with $\sin\psi\in(0,\beta)$. From $|v_{i+1}-w|\le\Dt$ we get $\sin\psi\ge1/(2\Dt)=\sin\psi^*$, where $v^*-v_i=(\cos\psi^*,\sin\psi^*)$. Suppose $\psi>\psi^*$. Write $u_{v_{i+1}}=(\sin\phi,\cos\phi)$; the T-edge is good, so $|\phi|\le\beta$. Since $v_i\in K$, $(v_i-v_{i+1})\cdot u_{v_{i+1}}=-\sin(\psi+\phi)\ge0$, so $\phi\le-\psi<0$. Since $v^*\in K$ (\cref{lem:arcK}),
\[
(v^*-v_{i+1})\cdot u_{v_{i+1}}=(\cos\psi^*-\cos\psi)\sin\phi+(\sin\psi^*-\sin\psi)\cos\phi\ge0 .
\]
But both terms are negative. Hence $\psi=\psi^*$ and $v_{i+1}=v^*$.
\end{proof}

So the walk moves along $\gamma_{AB}$ in steps of exactly $\alpha_2$ and cannot pass $v_B$ (that would put a point at distance less than $1$ from $v_B$). It reaches $v_B$ only if $j\alpha_2=k\alpha_1$ for some integer $j\ge1$. A floating-point search (\texttt{r4\_regionIII\_checks.py}) found no violation of \cref{lem:rigid} in $2\cdot10^5$ samples.

\begin{lemma}[arithmetic]\label{lem:arith}
There are no integers $j,k\ge1$ with $j\alpha_2=k\alpha_1$ and~\eqref{eq:kalpha}.
\end{lemma}
\begin{proof}
Since $\Dt<\Delta$, $\alpha_2>\alpha_1$, so $j\alpha_2=k\alpha_1$ forces $k>j$, and $j(\alpha_2-\alpha_1)=(k-j)\alpha_1\ge\alpha_1$. Hence $k\alpha_1=j\alpha_2\ge\alpha_1\alpha_2/(\alpha_2-\alpha_1)=1/(1/\alpha_1-1/\alpha_2)$. Put $f(\rho):=1/(2\arcsin(1/(2\rho)))$. From $x\le\arcsin x\le x+x^3/(6(1-x^2))$ (the Taylor coefficients of $\arcsin$ beyond the first are at most $1/6$) we get $\rho-1/(23\rho)\le f(\rho)\le\rho$ for $\rho\ge10^4$. So $1/\alpha_1-1/\alpha_2=f(\Delta)-f(\Dt)\le t+1/(23\Dt)<0.8662$, and $k\alpha_1>1/0.8662>1.154>\pi/3$, contradicting~\eqref{eq:kalpha}. (The inequality $t+1/(23\Dt)<3/\pi$ is checked in exact rational arithmetic in \texttt{r4\_regionIII\_checks.py}; a floating-point search found $\min_{j,k}|j\alpha_2-k\alpha_1|\,\Dt=1.7\cdot10^{-6}$, nonzero as it must be.)
\end{proof}

So for every pair $A,B$ of consecutive triangles with good bases in the same $Z$, the walk of \cref{lem:rigid} stops at a K-row vertex $b_{AB}$ on the half-open arc $[v_A,v_B)$ of $\gamma_{AB}$ with no pure right T-neighbour.

\begin{proof}[Proof of \cref{prop:tauexact}]
\emph{The half-open arcs $[v_A,v_B)$ are pairwise disjoint.} Arcs from the same $Z$ are consecutive intervals of one circle. For two components $Z\ne Z'$ with the same centre $w$, the angular ranges about $w$ (closed intervals of length less than $\pi/3$) are disjoint: a point of $Z'$ in the range of $Z$ either is a point of $Z$, which is impossible, or is at angle less than $\alpha_1$, hence at distance less than $1$, from a point of $Z$. So neither range contains an endpoint of the other, and the two intervals are disjoint. The arcs $[v_A,v_B)$ lie in the ranges of their components, so they are disjoint. For arcs $\gamma\subseteq\Cc(w,\Dt)$ and $\gamma'\subseteq\Cc(w',\Dt)$ of $\partial K$ with $w\ne w'$ that share a point $x$: if $x$ were in the relative interior of $\gamma$, then near $x$ the curve $\partial K$ coincides with $\gamma$, and the nondegenerate arc $\gamma'$ contains a piece of $\partial K$ on one side of $x$, which would then lie on both circles, forcing $w=w'$. So $x$ is an endpoint of both arcs; for half-open arcs this means $x=v_A=v_{A'}$, and since each K-vertex is the apex of at most one triangle, the two arcs come from the same triangle, so their centres agree, a contradiction.
Hence the vertices $b_{AB}$ are distinct, and $\mathrm{br}_K\ge\sum_Z(\nu_Z-1)$, where $\nu_Z$ is the number of triangles with good base in $Z$. Each $Z$ with $\nu_Z\ge1$ is a path (\cref{lem:onecircle}), and its right end has no pure right T-neighbour; distinct components have distinct right ends. So $\mathrm{br}_D\ge\#\{Z:\nu_Z\ge1\}$. Adding, $\mathrm{br}_K+\mathrm{br}_D\ge\sum_Z\nu_Z\ge q_2-C_{\mathrm{bad}}$, since every K-vertex with two rungs is the apex of a triangle, distinct apexes have distinct bases, and at most $C_{\mathrm{bad}}$ bases are bad. With~\eqref{eq:breaks} this proves the proposition.
\end{proof}

This completes the proof of \cref{prop:regionIII2}, and hence of \cref{prop:eS}. At $\tau=\sqrt3/2$ we did not use (R1) at D-vertices; it can fail there, since $2\sigma+2\beta>1$.

\section{Loss 2: points outside \texorpdfstring{$S$}{S}}\label{sec:loss2}

Recall from \cref{lem:identity} that each $y\in R$ contributes $(\deg y+m_y)/2$ to $\mult(\delta)$. In this section we prove that this is at most $5$, apart from $O(1)$ points, in every regime, and at most $4$ in Region~II. The section does not use the T-edge and rung analysis, and no hypothesis on $\tau$ enters except in \cref{thm:N2}.

\begin{lemma}[structure of $R$]\label{lem:Rstructure}
\begin{enumerate}[label=(\alph*)]
\item If $y\in R\cap D$, then $\deg y\le3$ and $\deg y+m_y\le6$.
\item $\partial K\subseteq\bigcup_{w\in X}\Cc(w,\Dt)$. Hence $S\setminus D=(X\setminus D)\cap\partial K$ and $R\setminus D\subseteq\operatorname{int}K$.
\item Every neighbour (point at distance $1$) of a point $y\in R\setminus D$ lies in $\operatorname{int}K$ (and then in $R\setminus D$), in $S\setminus D\subseteq\partial K$, or in $D$. Put $h:=\dist(y,\partial K)>0$. If $y$ has a neighbour in $S$, then $h\le1$, and if it has a neighbour in $D$, then $h\le1-t$.
\end{enumerate}
\end{lemma}
\begin{proof}
(a) \Cref{lem:identities}(a) holds for every point of $D$ (its proof uses only $|y-x_y|=\Delta$), so $X\setminus\{y\}$ lies in an open half-plane bounded by a line through $y$. Neighbours of $y$ are at least $60^\circ$ apart as seen from $y$, so there are at most $3$ of them.
(b) $X$ is finite, so a point of $K$ at distance less than $\Dt$ from every $w\in X$ is an interior point of $K$. A point $x\in(X\setminus D)\cap\partial K$ therefore lies on some $\Cc(w,\Dt)$ with $w\in X$, and $w\ne x$, so $x\in S$. Together with $X\setminus D\subseteq K$ and $S\setminus D\subseteq\partial K$ (\cref{lem:ballpoly}) this gives both statements.
(c) A neighbour not in $D$ lies in $K$ (\cref{lem:ballpoly}(a)), in $\partial K$ exactly when it is in $S$ by~(b). A neighbour $q\in S$ is not in $\operatorname{int}K$: either $q\in S\setminus D\subseteq\partial K$, or $q\in D$, and $\dist(q,K)\ge t$ by \cref{lem:ballpoly}(c). The segment $[y,q]$ leaves $K$ at a point $b\in\partial K$, so $h\le|yb|\le1$, and if $q\in D$ even $h\le1-|bq|\le1-t$.
\end{proof}

By~(a) and~(c), points $y\in R$ with $y\in D$ or $h>1$ have $\deg y+m_y\le6$. From now on let $y\in R\setminus D$ with $h\le1$, let $p_0\in\partial K$ be a nearest boundary point, $n_0:=(p_0-y)/h$ (an outer normal of $K$ at $p_0$), and $H_0:=\{x:(x-p_0)\cdot n_0\le0\}\supseteq K$, with boundary line $\ell_0$.

\emph{Small $K$.} If the perimeter of $K$ is less than $464$, then $\operatorname{diam}K<232$ and, since the discs of radius $\frac12$ about points of $X\setminus D\subseteq K$ are disjoint, $|X\setminus D|\le465^2$. Otherwise let $\gamma_y$ be the closed arc of $\partial K$ of points at arc-length distance at most $232$ from $p_0$.

\begin{lemma}[thin $K$]\label{lem:thin}
Suppose the perimeter of $K$ is at least $464$. If some $b\in\partial K$ with $|yb|\le1$ is not in $\gamma_y$, then $K$ lies in a rectangle with sides less than $400$ and $9$; in particular $|X\setminus D|<5106$.
\end{lemma}
\begin{proof}
Both arcs of $\partial K$ from $p_0$ to $b$ have length greater than $232$. Let $\nu_b$ be an outer normal of $K$ at $b$, $H_b:=\{x:(x-b)\cdot\nu_b\le0\}\supseteq K$ and $\ell_b:=\partial H_b$. We have $\dist(y,\ell_0)=h\le1$ and $\dist(y,\ell_b)\le|yb|\le1$.

\emph{Case $\nu_b=n_0$.} Since $p_0\in H_b$ and $b\in H_0$, $(p_0-b)\cdot n_0\le0\le(p_0-b)\cdot n_0$, so $\ell_0=\ell_b$. Then $[p_0,b]\subseteq K\cap\ell_0\subseteq\partial K$ is a boundary arc from $p_0$ to $b$ of length $|p_0b|\le2<232$, a contradiction.

\emph{Case $\nu_b\ne\pm n_0$.} The lines $\ell_0,\ell_b$ meet at a point $A$, and $\Omega:=H_0\cap H_b$ is a closed cone with apex $A$ and opening $\theta=\pi-\angle(n_0,\nu_b)\in(0,\pi)$. The points of $\Omega$ at distance at most $1$ from both bounding lines form a rhombus whose far vertex lies on the bisector at distance $1/\sin(\theta/2)$ from $A$, so $|yA|\le1/\sin(\theta/2)$. The point $p_0$ lies on the boundary ray of $\Omega$ on $\ell_0$ and $b$ on the one on $\ell_b$, and $|Ap_0|,|Ab|\le1+1/\sin(\theta/2)$. If $[p_0,b]\subseteq\partial K$ we conclude as in the first case. Otherwise the line $p_0b$ meets $\partial K$ only at $p_0$ and $b$; let $H^+$ be the closed half-plane bounded by it that contains $A$. Then $\partial K\cap H^+$ is one of the two arcs from $p_0$ to $b$, and $K\cap H^+\subseteq\Omega\cap H^+$, which is the triangle $Ap_0b$. By monotonicity of perimeter for convex sets, $\operatorname{length}(\partial K\cap H^+)\le|Ap_0|+|Ab|\le2+2/\sin(\theta/2)$. If $\theta\ge1^\circ$ this is less than $231.2<232$, a contradiction. So $\theta<1^\circ$.

\emph{Case $\theta<1^\circ$ or $\nu_b=-n_0$ (fatness step).} In the second case $\Omega$ is a strip of width at most $2$ containing $y$, and we put $\theta:=0$. Let $e$ be the direction of the bisector of $\Omega$ (of the strip if $\theta=0$), oriented so that $\Omega$ opens towards increasing $e$-coordinate, and $e^\perp$ a unit normal to $e$; call the $e$-coordinate the \emph{level} and the $e^\perp$-coordinate the \emph{height}. The cross-section of $\Omega$ at level $\lambda$ has length $w(\lambda_y)+2\tan(\theta/2)(\lambda-\lambda_y)$, where $\lambda_y$ is the level of $y$ and $w(\lambda_y)\le2/\cos(\theta/2)=:w_y$ (at levels below $\lambda_y$ it is shorter). Let $L>0$ be the length of the range of levels of $K$ and $W$ the length of its range of heights. Then $W\le w_y+2\tan(\theta/2)L$ and $L\le\operatorname{diam}K\le2\Dt$.

Pick $a,a'\in K$ at the lowest and highest level, let $\lambda_m$ be the middle level, and let $p_{\mathrm{top}},p_{\mathrm{bot}}$ be the top and bottom of the section of $K$ at level $\lambda_m$, a segment of positive length. Let $\nu=\cos\varphi\,e^\perp+\sin\varphi\,e$ be an outer normal at $p_{\mathrm{top}}$. Writing $a-p_{\mathrm{top}}=-\frac L2e+v_ae^\perp$ and $a'-p_{\mathrm{top}}=\frac L2e+v_{a'}e^\perp$ with $|v_a|,|v_{a'}|\le W$, the conditions $\nu\cdot(a-p_{\mathrm{top}})\le0$ and $\nu\cdot(a'-p_{\mathrm{top}})\le0$ read
\[
-\tfrac L2\sin\varphi+v_a\cos\varphi\le0,\qquad \tfrac L2\sin\varphi+v_{a'}\cos\varphi\le0 .
\]
If $\cos\varphi=0$ these two inequalities contradict each other, as $L>0$. If $\cos\varphi<0$, adding them gives $v_a+v_{a'}\ge0$; but the midpoint of $aa'$ lies in $K$ at level $\lambda_m$, so $v_a+v_{a'}\le0$. Hence $v_a+v_{a'}=0$, the midpoint of $aa'$ is $p_{\mathrm{top}}$, and $a,a'$ lie on the supporting line at $p_{\mathrm{top}}$ with normal $\nu$; then $\nu\cdot(p_{\mathrm{bot}}-p_{\mathrm{top}})>0$, contradicting $p_{\mathrm{bot}}\in K$. So $\cos\varphi>0$, and the inequalities give $|\tan\varphi|\le2W/L$. The same holds at $p_{\mathrm{bot}}$ with $-e^\perp$ in place of $e^\perp$.

By \cref{lem:Rstructure}(b), $p_{\mathrm{top}}\in\Cc(w_{\mathrm{top}},\Dt)$ for some $w_{\mathrm{top}}\in X$, and since $K\subseteq\Dd(w_{\mathrm{top}},\Dt)$, $(p_{\mathrm{top}}-w_{\mathrm{top}})/\Dt$ is an outer normal at $p_{\mathrm{top}}$; similarly for $p_{\mathrm{bot}}$. Applying the above to these normals,
\[
(w_{\mathrm{bot}}-w_{\mathrm{top}})\cdot e^\perp\ \ge\ -W+2\Dt\cos\varphi^*,\qquad\tan\varphi^*:=2W/L,
\]
while $|w_{\mathrm{bot}}-w_{\mathrm{top}}|\le\Delta\le1.94\,\Dt$. Suppose $L\ge400$. Then $\tan\varphi^*\le2w_y/400+4\tan(\theta/2)$ and $W\le w_y+4\Dt\tan(\theta/2)$, and so $(2\Dt\cos\varphi^*-W)/\Dt\ge2\cos\varphi^*-4\tan(\theta/2)-w_y/10^4>1.94$ for $\theta<1^\circ$; this contradicts the previous display. Hence $L<400$ and $W\le w_y+800\tan(0.5^\circ)<9$. So $K$ lies in a $400\times9$ rectangle, and the disjoint discs of radius $\frac12$ about the points of $X\setminus D$ lie in a $401\times10$ rectangle, so $|X\setminus D|<4010/(\pi/4)<5106$.
\end{proof}

\begin{lemma}[big turning]\label{lem:bigturn}
Suppose the perimeter of $K$ is at least $464$ and let $\eps_0>0$. Fewer than $2\pi\cdot468^2/\eps_0$ points $y\in R\setminus D$ with $h\le1$ have $\turn(\gamma_y)>\eps_0$. For $\eps_0=1/50$ this is fewer than $6.9\cdot10^7$.
\end{lemma}
\begin{proof}
If $b\in\gamma_y$ then $|b-p_0|\le232$, so $|b-y|\le233$. The points $y$ are at least $1$ apart, so for fixed $b$ at most $(2\cdot233.5)^2\le468^2$ of them qualify (disjoint discs of radius $\frac12$ in $\Dd(b,233.5)$). So every boundary point lies in at most $468^2$ of the arcs $\gamma_y$ (indexed by $y$), and \cref{lem:multiplicity} gives the bound.
\end{proof}

\begin{lemma}[flat case]\label{lem:flat}
Let $0<\eps_0\le1/50$ and $\eps:=\tan\eps_0$. Suppose the perimeter of $K$ is at least $464$, $\partial K\cap\Dd(y,1)\subseteq\gamma_y$ and $\turn(\gamma_y)\le\eps_0$. Use coordinates $(\xi,\zeta)$ with origin $p_0$ and $\zeta$-axis $n_0$, so $y=(0,-h)$. Then there is a concave function $f$ on $[-231,231]$ with $f(0)=0$, $f\le0$ and $|f(\xi)|\le\eps|\xi|$, whose graph is contained in $\gamma_y$, and $U:=\Dd(y,1)\cap\{\zeta<f(\xi)\}\subseteq\operatorname{int}K$. Write a neighbour of $y$ as $q=y+(\sin\psi,\cos\psi)$ with $\psi\in(-\pi,\pi]$. Then:
\begin{enumerate}[label=(F\arabic*)]
\item every neighbour $q\notin\operatorname{int}K$ (in particular every neighbour in $S$ and every neighbour in $D$) has $\cos\psi\ge h+f(\sin\psi)\ge h-\eps|\sin\psi|$; hence $|\psi|<\pi/2+\eps_0$;
\item every neighbour in $S\setminus D$ has $\cos\psi=h+f(\sin\psi)\in[h-\eps,h]$;
\item every neighbour in $D$ (in $S$ or not) has $\cos\psi\ge h+t+f(\sin\psi)\ge h+t-\eps$;
\item every neighbour in $\operatorname{int}K$ has $\cos\psi<h$.
\end{enumerate}
\end{lemma}
\begin{proof}
The normal set $N(\gamma_y)$ is an arc of $S^1$ of length at most $\eps_0$ containing $n_0$. Parametrise $\gamma_y$ by arc length $\lambda\in[-232,232]$ with $\gamma(0)=p_0$. Every one-sided tangent is perpendicular to a normal in $N(\gamma_y)$, so after orienting, the $\xi$-component of $\gamma'$ is at least $\cos\eps_0>0$. So $\xi(\lambda)$ is strictly increasing, $|\xi(\pm232)|\ge232\cos\eps_0>231$, and $\gamma_y$ contains the graph of a function $f$ on $[-231,231]$ with $|f'|\le\tan\eps_0=\eps$ almost everywhere and $f(0)=0$, so $|f(\xi)|\le\eps|\xi|$. Also $f\le0$ as $K\subseteq H_0$. At each graph point the outer normals have $\zeta$-component at least $\cos\eps_0>0$, and $K$ lies in the corresponding supporting half-plane, i.e.\ below the supporting line; hence $f$ is concave, and points directly above a graph point are not in $K$. The set $U$ is convex (the $\xi$-range of $\Dd(y,1)$ is $[-1,1]$), contains $y$ (as $-h<0=f(0)$), and misses $\partial K$: a point of $\partial K\cap\Dd(y,1)$ lies in $\gamma_y$, and the points of $\gamma_y$ with $|\xi|\le231$ are graph points, where $\zeta=f(\xi)$. A connected set that meets $\operatorname{int}K$ and misses $\partial K$ lies in $\operatorname{int}K$. So $U\subseteq\operatorname{int}K$.
(F1) A neighbour $q\notin\operatorname{int}K$ is not in $U$, so $\cos\psi-h\ge f(\sin\psi)$. Hence $\cos\psi>-\eps|\sin\psi|=-\tan\eps_0|\sin\psi|$, which excludes $|\psi|\ge\pi/2+\eps_0$. Neighbours in $S$ or in $D$ are not in $\operatorname{int}K$ by \cref{lem:Rstructure}.
(F2) Such a neighbour lies in $\partial K\cap\Dd(y,1)$, hence on the graph.
(F3) The graph point $g=(\sin\psi,f(\sin\psi))\in K$ lies directly below $q$ by (F1), so $t\le\dist(q,K)\le|q-g|=\cos\psi-h-f(\sin\psi)$ by \cref{lem:ballpoly}(c).
(F4) If $\cos\psi-h\ge f(\sin\psi)$, then $q$ is on the graph (so in $\partial K$) or directly above a graph point (so outside $K$). Hence a neighbour in $\operatorname{int}K$ has $\cos\psi-h<f(\sin\psi)\le0$.
\end{proof}

\begin{theorem}[N1]\label{thm:N1}
In the nondegenerate regime with $\Dt\ge10^4$, and for every value of $\tau$, all but fewer than $7.1\cdot10^7$ points $y\in R$ satisfy $m_y\le4$, and hence $\deg y+m_y\le10$.
\end{theorem}
\begin{proof}
Points in $R\cap D$ and points with $h>1$ have $m_y\le3$ and $m_y=0$ respectively (\cref{lem:Rstructure}). If the perimeter of $K$ is less than $464$, there are at most $465^2$ points in $R\setminus D$. Otherwise let $y\in R\setminus D$ with $h\le1$. If some point of $\partial K\cap\Dd(y,1)$ is not in $\gamma_y$, then $|X\setminus D|<5106$ by \cref{lem:thin}. If $\turn(\gamma_y)>1/50$, $y$ is one of fewer than $6.9\cdot10^7$ points (\cref{lem:bigturn}). In the remaining case \cref{lem:flat} applies with $\eps_0=1/50$, and by (F1) every neighbour of $y$ in $S$ has $|\psi|<\pi/2+\eps_0$, an open arc of length less than $4\pi/3$. Neighbours of $y$ are at least $60^\circ$ apart as seen from $y$, and an open arc of length less than $240^\circ$ contains at most $4$ such points. So $m_y\le4$. The exceptional points number less than $6.9\cdot10^7+5106+465^2<7.1\cdot10^7$.
\end{proof}

No hypothesis on $\tau$ and nothing about rungs enters \cref{thm:N1}; it holds in Regions~I, II and~III, at $\tau=\sqrt3/2$, for $\tau\le0$ and for $\tau>1$.

\begin{theorem}[N2]\label{thm:N2}
Let $0<\eps_0\le1/50$, $\eps:=\tan\eps_0$ (so $\eps<0.0201$), and suppose $t>\sqrt3/2+\eps$ (so $t>\frac12+2\eps$). In the nondegenerate regime with $\Dt\ge10^4$, all but fewer than $2\pi\cdot468^2/\eps_0+5106+465^2$ points $y\in R$ satisfy $m_y\le3$, and $m_y\le2$ if $\deg y=6$; hence $\deg y+m_y\le8$. In Region~II this holds with $\eps_0=1/50$, and then the exceptional set has fewer than $7.1\cdot10^7$ points.
\end{theorem}
\begin{proof}
The exceptions are as in \cref{thm:N1}; so let $y$ be in the case of \cref{lem:flat} (the other points have $\deg y+m_y\le6$).
\emph{(i) At most one neighbour in $D$.} By (F3) a neighbour in $D$ has $\cos\psi\ge t-\eps>\sqrt3/2$, i.e.\ $|\psi|<30^\circ$, and two neighbours at least $60^\circ$ apart cannot both lie in this open arc of $60^\circ$.
\emph{(ii) At most one neighbour in $S\setminus D$ with $\psi\in[0,\pi)$, and at most one with $\psi\in(-\pi,0]$.} Suppose $0\le\psi_1<\psi_2$, $\psi_2-\psi_1\ge60^\circ$, both in $S\setminus D$. By (F1), $\psi_2<90^\circ+\eps_0<91.2^\circ$ (recall $\eps_0\le1/50<1.2^\circ$), so $(\psi_1+\psi_2)/2\in[30^\circ,91.2^\circ]$ and $(\psi_2-\psi_1)/2\in[30^\circ,45.6^\circ]$. Then $\cos\psi_1-\cos\psi_2=2\sin\frac{\psi_1+\psi_2}2\sin\frac{\psi_2-\psi_1}2\ge\frac12>\eps$, contradicting (F2).
So $m_y\le3$, and if $m_y=3$ the neighbours in $S$ are one point of $D$ at some $\psi_D$ with $|\psi_D|<30^\circ$ and one point of $S\setminus D$ on each side.
\emph{(iii) $\deg y=6$ and $m_y=3$ is impossible.} Six neighbours pairwise at least $60^\circ$ apart are at the directions $\psi_D+60^\circ k$. The two neighbours in $S\setminus D$ are at two of the five directions $\psi_D\pm60^\circ$, $\psi_D\pm120^\circ$, $\psi_D+180^\circ$, and by (F2) their cosines differ by at most $\eps$.
\begin{itemize}[leftmargin=2em]
\item $\psi_D+180^\circ$: $\cos=-\cos\psi_D<-\sqrt3/2<-\eps$, contradicting (F1).
\item $\{\psi_D+60^\circ,\psi_D-60^\circ\}$: $|\cos(\psi_D+60^\circ)-\cos(\psi_D-60^\circ)|=\sqrt3|\sin\psi_D|\le\eps$, so $\cos(\psi_D\pm60^\circ)\ge\frac12\cos\psi_D-\frac{\sqrt3}2|\sin\psi_D|\ge\frac12-\eps$. By (F2), $h\ge\frac12-\eps$, and then (F3) gives $\cos\psi_D\ge h+t-\eps\ge\frac12+t-2\eps>1$, a contradiction.
\item $\{\psi_D+120^\circ,\psi_D-60^\circ\}$ (or its mirror image): $\cos(\psi_D-60^\circ)=-\cos(\psi_D+120^\circ)$, so both lie in an interval of length $\eps$ only if $|\cos(\psi_D+120^\circ)|\le\eps/2$; then $h\le3\eps/2<0.031$ by (F2), and since $\psi_D+120^\circ\in(90^\circ,150^\circ)$, $\psi_D+120^\circ\le90^\circ+\arcsin(\eps/2)$, so $\psi_D+60^\circ\in(30^\circ,31.2^\circ]$. The neighbour at $\psi_D+60^\circ$ is not in $S$ (as $m_y=3$), and $\cos(\psi_D+60^\circ)>\frac12>h$, so by (F4) it is not in $\operatorname{int}K$; by \cref{lem:Rstructure}(c) it is in $D$. But by (F3) every neighbour in $D$ has $|\psi|<30^\circ$. Contradiction.
\item $\{\psi_D+120^\circ,\psi_D-120^\circ\}$: both directions have $|\psi|>90^\circ$, so their cosines are negative and $h\le\eps$ by (F2). The neighbours at $\psi_D\pm60^\circ$ are not in $S$; one of them has $|\psi|\le60^\circ$, so $\cos\psi\ge\frac12>h$, and as before it lies in $D$, while $|\psi|>30^\circ$. Contradiction with (F3).
\item Two directions on the same side ($\{\psi_D+60^\circ,\psi_D+120^\circ\}$ or $\{\psi_D-60^\circ,\psi_D-120^\circ\}$) are excluded by~(ii).
\end{itemize}
Hence $\deg y=6$ implies $m_y\le2$, and $\deg y\le5$ gives $\deg y+m_y\le8$.
In Region~II, if $t\le1$ then $t\ge\tau\ge\sqrt3/2+3\beta$ (\cref{sec:setup}), and otherwise $t>1$; in both cases $t>\sqrt3/2+0.0201>\sqrt3/2+\tan(1/50)$. The count is $2\pi\cdot468^2\cdot50+5106+465^2<7.1\cdot10^7$.
\end{proof}

The constants in \cref{lem:thin,lem:bigturn} and the linear programmes are checked in exact rational arithmetic, with rigorous rational bounds for $\sin$, $\cos$ and $\tan$, in \texttt{r4\_nb\_exact.py} and \texttt{r4\_nb\_referee\_checks.py} (apex arc $2+2/\sin(0.5^\circ)<231.19$, the fatness margin, $W<9$, $M\le468^2$, the exceptional count).

The hypothesis of \cref{thm:N2} also holds in Region~III when $\tau\ge\sqrt3/2+\eta$ for a fixed $0<\eta\le1/50$, with $\eps_0=\arctan(\eta/2)$, but the exceptional count grows like $1/\eta$, so it is not uniform. For $\tau\le\sqrt3/2$ only \cref{thm:N1} is available. Locally it cannot be improved there:

\begin{remark}[local sharpness of \cref{thm:N1}]\label{rem:sharpN1}
Consider a single circular arc of $\partial K$: $K=\Dd(w,\Dt)$ near $y$, with $|y-w|=\Dt-h$. A neighbour at angle $\psi$ from the outward radial direction lies on $\partial K$ if and only if $\cos\psi=c_K(h):=h+(h^2-1)/(2(\Dt-h))$, and at distance at least $t$ from $K$ if and only if $\cos\psi\ge c_D(h):=h+t+((h+t)^2-1)/(2(\Dt-h))$. One checks $c_D(0)=\tau$, and $c_D$ is increasing in $h$. So two neighbours $60^\circ$ apart that are both at distance at least $t$ from $K$ need $\tau<\sqrt3/2$. For $\tau<\sqrt3/2$ this is attained: take $w=0$, $\Dt=10^4$, a chord $[p_-,p_+]$ of length $2$ of $\Cc(w,\Dt)$, $y$ its midpoint, $e:=(p_+-p_-)/2$, $\nu:=y/|y|$, $z_\pm:=y\pm e/2+\frac{\sqrt3}2\nu$ and $a_\pm:=y\pm e/2-\frac{\sqrt3}2\nu$. Then $|yp_\pm|=|yz_\pm|=|ya_\pm|=1$, all pairwise distances are at least $1$, $z_\pm\in\Cc(w,\Dt+t)$ with $\dist(z_\pm,\Dd(w,\Dt))=t$ for $t=0.86598790\ldots$ ($\tau=\sqrt3/2-5.0\cdot10^{-5}$), and $a_\pm\in\operatorname{int}\Dd(w,\Dt)$. This is verified in $60$-digit decimal arithmetic in \texttt{r4\_nb\_exact.py}. So $y$ satisfies every local constraint used above with $\deg y=6$, two neighbours on $\partial K$ and two neighbours that may lie in $D$; if the latter lie in $S$, then $\deg y+m_y=10$. In this one-centre example the points $z_\pm$ have no $\Dt$-partner, and whether $\Theta(n)$ such points $y$, with the $z_\pm$ in $S$, can occur in a set $X$ is open. Hence local arguments of the kind used here cannot give a bound below $10$ for $\tau<\sqrt3/2$.
\end{remark}

\section{Proof of Theorem 1}\label{sec:proof15}

\begin{table}[ht]
\centering
\renewcommand{\arraystretch}{1.2}
\begin{tabular}{@{}llll@{}}
\toprule
regime ($\delta=1$, $\beta=\frac1{100}$) & $e(S)\le$ & $\deg y+m_y$ & $\min\{\mult(\Dt),\mult(\delta)\}\le$\\
\midrule
degenerate: $\Delta>1.94\Dt$, $n\ge4$ & --- & --- & $n+1$ (\cref{prop:degenerate})\\
nondegenerate, $\Dt<10^4$ & --- & --- & $3n<3N_0$ (\cref{sec:reductions})\\
I: $\tau\le\sqrt3/2-2\beta$ & $s+C$ & $\le10$ (\cref{thm:N1}) & $\frac{15}{11}n+C$\\
III: $\sqrt3/2-2\beta<\tau<\sqrt3/2+3\beta$ & $s+p+C$ & $\le10$ (\cref{thm:N1}) & $\frac{15}{11}n+C$\\
II: $\sqrt3/2+3\beta\le\tau\le1$ & $s+p+C$ & $\le8$ (\cref{thm:N2}) & $\frac43n+C$\\
II: $\tau>1$ & $s+C$ & $\le8$ (\cref{thm:N2}) & $\frac43n+C$\\
\bottomrule
\end{tabular}
\caption{The regimes in the proof of Theorem 1. The bounds on $e(S)$ are \cref{prop:eS}; the bounds on $\deg y+m_y$ hold for $y\in R$ outside a set of fewer than $7.1\cdot10^7$ points; here $C<2\cdot10^8$.}
\end{table}

In the nondegenerate regime with $\Dt\ge10^4$, \cref{prop:eS} gives $e(S)\le s+p+C_1$ with $C_1<3.7\cdot10^7$ in every $\tau$-regime (and $e(S)\le s+C_{\mathrm{bad}}$ in Region~I and for $\tau>1$). \Cref{thm:N1} gives $\deg y+m_y\le10$ for all but $C'<7.1\cdot10^7$ points of $R$, so \cref{lem:lp3}(a) gives $\min\{\mult(\Dt),\mult(\delta)\}\le\frac{15}{11}n+C_1+C'$, with $C_1+C'<1.1\cdot10^8$. In Region~II, \cref{thm:N2} gives $\deg y+m_y\le8$ for all but $C'<7.1\cdot10^7$ points, and \cref{lem:lp3}(b) gives $\frac43n+C_1+2C'$, with $C_1+2C'<1.8\cdot10^8$. In the degenerate regime $\mult(\Dt)\le n+1\le\frac{15}{11}n+1$ (\cref{prop:degenerate}); if $\Dt<10^4$ then $n<N_0$ and $\mult(\delta)\le3n<3N_0$ (\cref{sec:reductions}). The $\tau$-intervals partition $\R$, and the cases $n\le3$ are trivial. Hence Theorem~1 holds with $C_0=3N_0+2\cdot10^8$, including the refinement $\frac43n+C_0$ in Region~II. \qed

\begin{remark}\label{rem:ienjoy}
The bound $15/11$ is consistent with the lower bound $9/7$ of~\cite{ienjoymath}. That construction lies in Region~I (there $t/\delta\to0$); it has $e(S)=|S|$ and $Q=\emptyset$, and its points in $R$ have $m_y=0$ and $\deg y+m_y\le6$ (checked for $n=60,120,250$ in \texttt{r4\_nb\_ring16\_check.py}). With $\deg y+m_y\le6$ and $e(S)\le s+O(1)$, \cref{rem:lpgeneral} gives exactly $9/7$. The linear programme behind $15/11$ is extremal for $s=\frac{10}{11}n$ with points of $R$ having $\deg y+m_y=10$.
\end{remark}

\section{A layer bound and a new class of sets}\label{sec:layers}

In this section $X$ is an arbitrary finite planar set with at least two distinct distances, $\Delta>\Dt$ are its two largest distances, and we scale so that $\Dt=1$. Every distance of $X$ is at most $1$ or equal to $\Delta$. $L_1$ is the set of extreme points of $X$ (vertices of $\conv X$), $L_2$ is the set of extreme points of $X\setminus L_1$, and $m_3:=n-|L_1|-|L_2|$. Let $G$ be the graph on $L_1\cup L_2$ whose edges are the $\Dt$-pairs, let $G'$ be its $2$-core (repeatedly delete vertices of degree less than $2$), and write $L_i':=L_i\cap V(G')$ and $\deg'=\deg_{G'}$. For $p\in X$ let $h(p)$ be the number of points of $X$ at distance $\Dt$ from $p$.

\begin{lemma}[location]\label{lem:location}
\begin{enumerate}[label=(\roman*)]
\item Both endpoints of every $\Delta$-pair lie in $L_1$.
\item Every $\Dt$-pair has an endpoint in $L_1$ and both endpoints in $L_1\cup L_2$.
\end{enumerate}
In particular $G$ contains every $\Dt$-pair of $X$.
\end{lemma}
\begin{proof}
(i) If $|pq|=\Delta$ then $X\subseteq\Dd(q,\Delta)$ and $p$ lies on its boundary circle, so $p$ is an extreme point of $X$. (ii) Let $|pq|=1$. If $p\notin L_1$, then $p$ has no $\Delta$-partner by (i), so $X\subseteq\Dd(p,1)$ and $q$, lying on the boundary circle, is in $L_1$. So some endpoint, say $p$, lies in $L_1$. Suppose $q\notin L_1$. Every point $x$ with $|px|>1$ has $|px|=\Delta$ and so lies in $L_1$. Hence $X\setminus L_1\subseteq\Dd(p,1)$ with $q$ on the boundary circle, so $q$ is extreme in $X\setminus L_1$, i.e.\ $q\in L_2$.
\end{proof}

\begin{proposition}[core degree of $L_2$]\label{prop:core2}
Every $q\in L_2'$ has $\deg'(q)=2$.
\end{proposition}
\begin{proof}
Put $q=0$. Since $q$ is not extreme it has no $\Delta$-partner (\cref{lem:location}(i)), so $X\subseteq\Dd(0,1)$ and the $G$-neighbours of $q$ lie on the unit circle. Since $q\in V(G')$, $\deg'(q)\ge2$. Suppose $v_1,v_2,v_3$ are three distinct $G'$-neighbours of $q$. Each $|v_iv_j|$ is either at most $1$ or equal to $\Delta$. The argument follows the three cases of Vesztergombi's inner-point analysis~\cite[pp.~192--194]{Ve87}, with one change: we delete no points, and the contradiction in Cases~1 and~2 is with membership in the $2$-core, via $\deg_G\ge\deg'$.

\emph{Case 3: at most one pair $v_iv_j$ at distance $\Delta$.} Label so that $|v_1v_2|,|v_2v_3|\le1$. Then $v_2$ lies on the shorter arc between $v_1$ and $v_3$, and in polar angle we may take $v_2$ at $0$, $v_1$ at $a$ and $v_3$ at $-b$ with $0<a,b\le60^\circ$. Since $q$ is not extreme, the closed half-plane $\{x:x\cdot v_2\le0\}$ contains a point $t\in X$, $t\ne q$. For $t\ne0$ with $t\cdot v_2\le0$ we get $|t-v_2|^2=|t|^2-2t\cdot v_2+1>1$, so $|tv_2|=\Delta$. By symmetry assume the polar angle $\varphi$ of $t$ lies in $[90^\circ,180^\circ]$. Then the angle between $t$ and $v_3$ is $\varphi+b\in(90^\circ,240^\circ]$, so $t\cdot v_3<0$ and $|tv_3|=\Delta$. Hence $t$ lies on the perpendicular bisector of $v_2v_3$, which passes through $0$, so $\varphi=180^\circ-b/2$. The angle between $t$ and $v_1$ is then $180^\circ-b/2-a\ge90^\circ$, so $|tv_1|^2\ge|t|^2+1>1$ and $|tv_1|=\Delta$. Now $v_1,v_2,v_3$ lie on both $\Cc(0,1)$ and $\Cc(t,\Delta)$, two distinct circles, which meet in at most two points. Contradiction. (This case uses only that $q$ is not extreme.)

\emph{Case 2: exactly two pairs at $\Delta$.} Say $|v_1v_2|=|v_1v_3|=\Delta$ and $|v_2v_3|\le1$. Then $v_1$ lies on the perpendicular bisector of $v_2v_3$, and we may take $v_1=(-1,0)$, $v_{2,3}=(\cos a,\pm\sin a)$ with $0<a\le30^\circ$ and $\Delta=2\cos(a/2)$. We claim $v_1$ has no $\Dt$-neighbour $r\ne q$; then $\deg_G(v_1)=1$, contradicting $v_1\in V(G')$. Such an $r=(x,y)$ satisfies $|r-v_1|=1$ (so $|r|^2=-2x$), $|r|\le1$, and $|r-v_i|\in[0,1]\cup\{\Delta\}$ for $i=2,3$; by symmetry $y\ge0$. Write $c=\cos a$, $s_a=\sin a$.
\begin{itemize}[leftmargin=2em]
\item \emph{Both $|r-v_2|,|r-v_3|\le1$.} Impossible: $|r-v_3|^2-1=-2x(1+c)+2ys_a\ge0$, with equality only if $r=q$.
\item \emph{$|r-v_2|=\Delta$} (the case $|r-v_3|=\Delta$ is symmetric). Using $|r|^2=-2x$, this condition is the linear equation $-2x(1+c)-2s_ay=1+2c$. Intersecting with the circle $|r-v_1|=1$ gives $x=\bigl(\pm s_a\sqrt{8c+7}-3(1+c)\bigr)/(4(1+c))$, so $|r|^2=\bigl(3(1+c)\mp s_a\sqrt{8c+7}\bigr)/(2(1+c))$. The root with $+s_a\sqrt{8c+7}$ in $|r|^2$ has $|r|^2\ge3/2$, outside the disc. For the other root, $|r|^2\le1\iff1+c\le s_a\sqrt{8c+7}\iff1+c\le(1-c)(8c+7)\iff c^2\le\frac34\iff a\ge30^\circ$. So only $a=30^\circ$ remains, and there $|r-v_3|^2=3+c-s_a\sqrt{8c+7}=2$, which lies strictly between $1$ and $\Delta^2=2+\sqrt3$. Forbidden.
\item \emph{Both equal $\Delta$.} Then $r$ lies on the $x$-axis, so $r\in\{(0,0),(-2,0)\}$, and $(-2,0)$ is outside the disc.
\end{itemize}

\emph{Case 1: all three pairs at $\Delta$.} Then $v_1,v_2,v_3$ form an equilateral triangle with circumcentre $q$, and $\Delta=\sqrt3$. Let $P\ne q$ be a further $\Dt$-neighbour of some $v_i$, so $|P-v_i|=1$ and $|P|\le1$. It cannot have both other distances at most $1$: $|P-v_i|=1$ gives $|P|^2=2P\cdot v_i$, the two conditions give $P\cdot(v_i-v_j)\le0$ for $j\ne i$, and adding and using $\sum_jv_j=0$ gives $3P\cdot v_i\le0$, so $P=q$. Hence $|P-v_j|=\sqrt3$ for some $j\ne i$. Solving exactly, there are six candidate points, two for each $v_i$, each with $|P|^2=(9-\sqrt{33})/6$ and at squared distance $(7-\sqrt{33})/2\approx0.63$ from the third vertex, so each is a $\Dt$-neighbour of exactly one $v_i$. Their pairwise squared distances are $(29-5\sqrt{33})/6\approx0.046$ for three disjoint pairs (each joining candidates of different vertices), and otherwise $4/3$, $(9-\sqrt{33})/2\approx1.628$ or $(17-\sqrt{33})/6\approx1.876$, all strictly inside the forbidden interval $(1,3)$. So a set of candidates that can coexist in $X$ has at most two elements, at most two of the $v_i$ have a second $\Dt$-neighbour, and some $v_i$ has $\deg_G(v_i)=1$, contradicting $v_i\in V(G')$.

With four or more neighbours apply the argument to any three. Hence $\deg'(q)\le2$, so $\deg'(q)=2$.
\end{proof}

\begin{lemma}[opposite-side lemma]\label{lem:M}
Let $|pa|=|pb|=|br|=1$ with $r\neq p$. Suppose $a$ lies strictly on one side of the line $pb$ and $r$ does not lie strictly on that side. Then $|ar|>1$, so $|ar|=\Delta$.
\end{lemma}
\begin{proof}
Put $p=0$, $b=(1,0)$, $a=(\cos A,-\sin A)$ with $0<A<\pi$, and $r=b+(\cos F,\sin F)$ with $0\le F\le\pi$; $F\ne\pi$ as $r\ne p$. Then $|r-a|^2-1=2(1+\cos F)(1-\cos A)+2\sin F\sin A$, where the first term is positive and the second nonnegative.
\end{proof}

\begin{lemma}\label{lem:F}
Let $p$ be an extreme point of $X$ with $\Dt$-neighbours $w_1,\dots,w_k$ in angular order. Then $\deg_G(w_i)=1$ for every $3\le i\le k-2$.
\end{lemma}
\begin{proof}
Since $p$ is extreme, some line through $p$ has $X\setminus\{p\}$ strictly on one side. So the directions of the $w_j$ from $p$ are distinct, lie in an open half-plane, and any two differ by less than $180^\circ$; for $j<i<l$, $w_j$ and $w_l$ lie strictly on opposite sides of the line $pw_i$. Suppose $r\ne p$ and $|rw_i|=1$. Put $p=0$, $w_i=(1,0)$. Then $w_{i-2},w_{i-1}$ lie strictly below the $x$-axis and $w_{i+1},w_{i+2}$ strictly above; by symmetry assume $r$ is not strictly below. Write $w_{i-2}=e(-\alpha_1)$, $w_{i-1}=e(-\alpha_2)$ with $0<\alpha_2<\alpha_1<180^\circ$. By \cref{lem:M}, $|rw_{i-2}|=|rw_{i-1}|=\Delta$, so $r$ lies on the perpendicular bisector of $w_{i-2}w_{i-1}$, the line through $0$ in direction $-\gamma$ with $\gamma=(\alpha_1+\alpha_2)/2\in(0^\circ,180^\circ)$. As $r\ne0$ is not below the axis, $r=\lambda e(180^\circ-\gamma)$ with $\lambda>0$. From $|r-w_i|=1$ we get $\lambda(\lambda+2\cos\gamma)=0$, so $\lambda=-2\cos\gamma$, whence $\gamma>90^\circ$ and $\alpha_1>90^\circ$. Then $|w_{i-2}w_i|=2\sin(\alpha_1/2)>\sqrt2>1$, so $|w_{i-2}w_i|=\Delta$. Let $\theta>0$ be the direction of $w_{i+1}$; then $\alpha_1+\theta<180^\circ$, and since the chord length $2\sin(x/2)$ is strictly increasing on $[0^\circ,180^\circ]$, $|w_{i-2}w_{i+1}|=2\sin((\alpha_1+\theta)/2)>2\sin(\alpha_1/2)=\Delta$, contradicting the maximality of $\Delta$.
\end{proof}

\begin{theorem}[degree at most $4$]\label{thm:deg4}
Every $p\in L_1'$ has $\deg'(p)\le4$.
\end{theorem}
\begin{proof}
Suppose $p\in L_1'$ has five $G'$-neighbours $u_1,\dots,u_5$ in angular order. Among the $G$-neighbours of $p$, $u_3$ has at least two on each side, so $u_3=w_i$ with $3\le i\le k-2$, and $\deg_G(u_3)=1$ by \cref{lem:F}. But $u_3\in V(G')$ forces $\deg_G(u_3)\ge2$.
\end{proof}

\begin{theorem}[layer bound]\label{thm:layer}
For every finite $X\subset\R^2$ with at least two distinct distances,
\[
\mult(\Dt)\ \le\ |L_1|+|L_2|+\tfrac12\sum_{p\in L_1'}\bigl(\deg'(p)-2\bigr).
\]
\end{theorem}
\begin{proof}
By \cref{lem:location}, $\mult(\Dt)=e(G)$. Each deletion made while pruning $G$ to its $2$-core removes at most one edge, so $e(G)\le e(G')+|V(G)\setminus V(G')|$. By \cref{prop:core2}, $e(G')=\frac12\sum_{v\in V(G')}\deg'(v)=|V(G')|+\frac12\sum_{p\in L_1'}(\deg'(p)-2)$. Add the two estimates and use $|V(G)|=|L_1|+|L_2|$.
\end{proof}

\begin{corollary}\label{cor:A1}
Assume \textup{(H3)}: no point of $L_1$ has four or more points of $X$ at distance $\Dt$. (This holds, for instance, if no four points of $X$ are concyclic.) Then
\[
\mult(\Dt)\le|L_1|+|L_2|+\tfrac12\#\{p\in L_1':\deg'(p)=3\}\le\tfrac32|L_1|+|L_2|.
\]
Consequently, under \textup{(H3)}, if $|L_1|\le2m_3$ then $\Delta$ and $\Dt$ are two distinct distances each occurring at least once and at most $n$ times.
\end{corollary}
\begin{proof}
Under (H3), $2\le\deg'(p)\le h(p)\le3$ for $p\in L_1'$. The condition $\frac32|L_1|+|L_2|\le n$ is equivalent to $|L_1|\le2m_3$, and $\mult(\Delta)\le n$ by~\cite{HP34}.
\end{proof}

\begin{corollary}[unconditional form]\label{cor:uncond}
For every finite $X$ with at least two distances,
$\mult(\Dt)\le|L_1|+|L_2|+\frac12\sum_{p\in L_1}(\min\{h(p),4\}-2)^+$. Hence $\mult(\Dt)\le n$ whenever $\sum_{p\in L_1}(\min\{h(p),4\}-2)^+\le2m_3$.
\end{corollary}
\begin{proof}
Combine \cref{thm:layer} with $2\le\deg'(p)\le\min\{h(p),4\}$ on $L_1'$ (\cref{thm:deg4}).
\end{proof}

\begin{remark}\label{rem:layers}
\begin{enumerate}[leftmargin=2em,label=(\alph*)]
\item \Cref{cor:A1} covers sets not covered by~\eqref{eq:cdl13}. For example, if $|L_1|=|L_2|=0.4n$ and $m_3=0.2n$, the three quantities in~\eqref{eq:cdl13} are $1.2n$, $1.33n$ and $1.2n$, while \cref{cor:A1} gives $\mult(\Dt)\le n$ (under (H3)). Of course~\eqref{eq:cdl13} needs no hypothesis like (H3).
\item Using $\deg'\le4$, \cref{thm:layer} recovers the bound $2|L_1|+|L_2|$ of~\cite[Theorem~1.3]{CDL25}. \Cref{prop:core2} and \cref{thm:deg4} are the inputs from~\cite{Ve96} (Propositions~3 and~5--8 there, as used in~\cite{CDL25}) needed for this; we prove them directly in the $2$-core, without deleting points. We have not checked the other inputs used in~\cite{CDL25} for the first two bounds in~\eqref{eq:cdl13}.
\item In \cref{cor:A1}, the bound $\frac32|L_1|+|L_2|$ is attained under (H3) alone by the following example, which we have checked numerically for $m=8,12,16$ (floating point): a regular $m$-gon with $m$ even, with a point at distance $\Dt$ from both endpoints of every other edge (CDL's ``servers''). These examples contain concyclic quadruples; for sets with no four concyclic points see \cref{q:genericity} below.
\end{enumerate}
\end{remark}

\section{Sets in convex position}\label{sec:convex}

A set $X$ is in \emph{convex position} if every point of $X$ is a vertex of $\conv X$; in particular no three points of $X$ are collinear. Throughout this section $|X|=n$, $m:=\lfloor n/2\rfloor$, $N:=\binom n2$, $D=D(X)$ is the number of distinct distances, and $k=k(X)$ is the number of rare distances ($1\le\mult(d)\le n$). A distance is \emph{frequent} if $\mult(d)\ge n+1$. $R_N$ denotes the vertex set of a regular $N$-gon. By Hopf--Pannwitz, $\Delta$ is always rare; Clemen, Dumitrescu and Liu proved $k\ge2$ for $n\ge5$~\cite[Theorem~1.2]{CDL25}. For $n=5$ this is best possible: $R_5$ has $k=2$.

\begin{theorem}\label{thm:convex-summary}
Let $X$ be a set of $n$ points in convex position with $6\le n\le13$. Then at least three distinct distances each occur at least once and at most $n$ times. The same holds for every even $n\ge6$.
\end{theorem}

The cases have different status:
\begin{itemize}[leftmargin=2em]
\item $n=6,7,8$: a corollary of published results (\cref{prop:lit678}).
\item even $n\ge10$: \cref{thm:even}, a parity-sharp version of the proof of~\cite[Theorem~1.2]{CDL25};
\item $n=9$: \cref{thm:nine}, from the classification of Erd\H{o}s and Fishburn~\cite{EF96};
\item $n=11,13$: \cref{thm:1113}, new and computer-assisted. The $n=7$ case is re-proved the same way, without the literature.
\end{itemize}

\subsection{Counting and the literature}

\begin{lemma}[parity-sharp counting]\label{lem:counting}
$(D-k)(n+1)+k\le N$. Consequently, if $n=2m+1$ and $D=m+K$ with $1\le K\le m+1$, then $k\ge K+1$; if $n=2m$ and $D=m+K$ with $1\le K<2m$, then $k\ge K+2$.
\end{lemma}
\begin{proof}
The first inequality holds because frequent distances contribute at least $n+1$ pairs each and rare ones at least $1$. It gives $k\ge(D(n+1)-N)/n$. For $n=2m+1$ the right side is $K+(m+K)/(2m+1)$, with $0<(m+K)/(2m+1)\le1$. For $n=2m$ it is $K+1+K/(2m)$, with $0<K/(2m)<1$.
\end{proof}

We use the following results from the literature.
\begin{itemize}[leftmargin=2em]
\item[(Alt)] Altman~\cite{Alt63,Alt72}: $D\ge m$; for odd $n$, $D=m$ only for $R_n$ (see also~\cite[Theorem~1]{Fi95}).
\item[(Fi)] Fishburn~\cite{Fi95}: for even $n\ge8$, $D=n/2$ if and only if $X$ is $R_n$ or $n$ vertices of $R_{n+1}$. For $n=6$ the hexagons with $D=3$ are $R_6$, $R_7$ minus a vertex, and one further hexagon~\cite[Theorem~2]{Fi95}, \cite{EF95b}, with multiplicity profiles $(6,6,3)$, $(5,5,5)$ and $(6,6,3)$ respectively (as recorded in~\cite[\S2]{CDL25}; we re-checked the profiles in exact arithmetic).
\item[(EF96)] Erd\H{o}s--Fishburn~\cite{EF96}: a convex nonagon has exactly $5$ distances if and only if its vertices are $9$ vertices of $R_{10}$ or of $R_{11}$.
\item[(r$_2$)] Erd\H{o}s--Fishburn~\cite[\S3]{EF95b}, quoting Fishburn's technical report~\cite{Fi92}: the largest possible second-largest multiplicity in a convex $n$-gon is $r_2^{(n)}=1,3,5,6,7,8$ for $n=3,\dots,8$.
\end{itemize}
A set on a circle has every distance class of maximum degree at most $2$, hence of size at most $n$; so for cocircular $X$ all distances are rare and $k=D$.

\begin{proposition}[from the literature]\label{prop:lit678}
For $n=6,7,8$ in convex position, $k\ge3$.
\end{proposition}
\begin{proof}
By (r$_2$), $r_2^{(n)}\le n$ for $n\le8$, so at most one distance is frequent and $k\ge D-1$. By (Alt), $D\ge3$ for $n=6,7$ and $D\ge4$ for $n=8$. So $k\ge3$ unless $n\in\{6,7\}$ and $D=3$. For $n=7$, $D=3$ forces $R_7$ (profile $(7,7,7)$, $k=3$); for $n=6$, (Fi) gives the profiles $(6,6,3)$ and $(5,5,5)$, so $k=3$.
\end{proof}
The value $r_2^{(7)}=7$, and the other values of $r_2^{(n)}$, rest on an unrefereed technical report~\cite{Fi92} quoted in a refereed paper~\cite{EF95b}; for $n=7$ we give an independent proof in \cref{thm:1113}.

\begin{theorem}[even $n$]\label{thm:even}
Let $X$ be in convex position with $n=2m\ge6$. Then $k\ge3$. If $n\ge8$ and $k\le3$, then $D=m+1$.
\end{theorem}
\begin{proof}
By (Alt), $D\ge m$. If $D\ge m+1$, \cref{lem:counting} gives $k\ge3$, and $k\ge4$ if $D\ge m+2$. If $D=m$ and $n=6$, (Fi) gives $k=3$. If $D=m$ and $n\ge8$, then $X$ is $R_n$, with profile $(n,\dots,n,n/2)$, or $n$ vertices of $R_{n+1}$, with profile $(n-1,\dots,n-1)$; in both cases $k=m\ge4$.
\end{proof}
\Cref{thm:even} is essentially the argument of~\cite[Theorem~1.2]{CDL25}, with the counting done with parity. It adds to the literature only for even $n\ge10$.

\begin{theorem}[$n=9$]\label{thm:nine}
Every set of $9$ points in convex position has $k\ge3$.
\end{theorem}
\begin{proof}
Here $m=4$. If $D=4$, then $X=R_9$ by (Alt) and $k=4$. If $D=5$, then $X$ is cocircular by (EF96) and $k=D=5$. If $D\ge6$, \cref{lem:counting} gives $k\ge3$.
\end{proof}

\subsection{\texorpdfstring{$n=7,11,13$}{n = 7, 11, 13}: an ordinal reduction and SAT certificates}\label{sec:sat}

Label $X=(p_0,\dots,p_{n-1})$ in convex cyclic order, let $d_1<\dots<d_D$ be the distinct distances, and colour each pair $\{i,j\}$ with $c(ij)=t$ where $|p_ip_j|=d_t$.

\begin{lemma}[quadrilateral inequality]\label{lem:Q}
For $i<j<k<l$ (in cyclic order), $|p_ip_k|+|p_jp_l|>|p_ip_j|+|p_kp_l|$ and $|p_ip_k|+|p_jp_l|>|p_jp_k|+|p_lp_i|$.
\end{lemma}
\begin{proof}
The diagonals of the convex quadrilateral $p_ip_jp_kp_l$ meet at a point $O$ that lies on neither of the sides $p_ip_j$, $p_kp_l$. So $|p_ip_k|+|p_jp_l|=(|p_iO|+|Op_j|)+(|p_kO|+|Op_l|)>|p_ip_j|+|p_kp_l|$, by the strict triangle inequality. The other inequality is the same.
\end{proof}

\emph{Ordinal form.} Let $x_1\le x_2$ be the two diagonal colours, and $y_1\le y_2$ the colours of one pair of opposite sides. Then it is \emph{not} the case that $x_1\le y_1$ and $x_2\le y_2$, since otherwise $|p_ip_k|+|p_jp_l|\le$ the sum of the two sides. We call a colouring of the pairs of an $n$-cycle's vertex set satisfying this for all $4$-tuples and both pairs of opposite sides a \emph{Q-colouring}. This ordinal use of the quadrilateral inequality goes back to Altman~\cite{Alt63}. Abstracting a convex distance problem to a finite search in this way is the method of Mori\'c and Pritchard~\cite{MP11}, whose Fact~3.1 is \cref{lem:Q}.

Some of our certificates also use a second valid constraint (\emph{Lemma T}): if $w\ne w'$ lie on the same open boundary arc between $p_i$ and $p_j$, then $(c(iw),c(jw))\ne(c(iw'),c(jw'))$, because a point is determined by its distances to $p_i,p_j$ and a side of the line $p_ip_j$.

\begin{lemma}[reduction to profiles]\label{lem:R}
If $k\le2$, then at least $D-2$ colour classes have size at least $n+1$, and $D\le D_{\max}(n)$, the largest $D$ with $(D-2)(n+1)+2\le N$. For $n=7,11,13,15$, $D_{\max}=4,6,7,8$.
\end{lemma}

Call a Q-colouring using exactly $D$ colours, all of them, with at least $F$ classes of size at least $n+1$, an \emph{$(n,D,F)$ ordinal counterexample}. Since \cref{lem:Q} (and Lemma T) are necessary for geometric realisability, the following suffices to prove $k\ge3$ for a given $n$: \emph{for every $2\le D\le D_{\max}(n)$ there is no $(n,D,D-2)$ ordinal counterexample.} (We check the smaller $D$ even where (Alt) would exclude them, so that no citation is needed; for $D\le4$ we check the stronger statement with $F=0$, i.e.\ no Q-colouring with $D$ colours at all.) No realisation step and no floating-point arithmetic enters.

\begin{theorem}\label{thm:1113}
Every set of $n$ points in convex position with $n\in\{7,11,13\}$ has at least three rare distances. For these $n$, the ordinal quadrilateral inequality (\cref{lem:Q}) alone implies the statement.
\end{theorem}
\begin{proof}[Proof (computer-assisted)]
By \cref{lem:R} it suffices that the following instances have no ordinal counterexample:
\begin{itemize}[leftmargin=2em]
\item $n=7$: $(7,2,0)$, $(7,3,1)$, $(7,4,2)$;
\item $n=11$: $(11,2,0)$, $(11,3,0)$, $(11,4,0)$, $(11,5,3)$, $(11,6,4)$;
\item $n=13$: $(13,2,0)$, $(13,3,0)$, $(13,4,0)$, $(13,5,0)$, $(13,6,4)$, $(13,7,5)$.
\end{itemize}
Each is encoded in CNF and shown unsatisfiable, with the refutation checked as described below and in \cref{app:verification}.
\end{proof}

\subsubsection*{Encoding}
Variables $x_{e,c}$ say that pair $e$ has colour $c$; there are exactly-one constraints per pair and at-least-one constraints per colour. In the \emph{direct} encoding of \cref{lem:Q}, each dominated $4$-tuple of colours is forbidden. In the \emph{order} encoding, auxiliary literals $y_{e,t}\Leftrightarrow c(e)\ge t$ are introduced and ``$c(a)\le t\le c(b)$ and $c(a')\le u\le c(b')$'' is forbidden for both diagonal/side matchings; this is equivalent, since sorted domination holds if and only if some matching dominates. Selector variables $f_c$ imply ``class $c$ has at least $n+1$ pairs'' through sequential counters, and $\sum_cf_c\ge F$.

\subsubsection*{Certification}
\begin{itemize}[leftmargin=2em]
\item \emph{DRAT/LRAT.} For every instance with $n=11$ and $n=13$ listed above, and for $(7,4,2)$, a DRUP proof produced by Glucose~4~\cite{glucose} was verified by \texttt{drat-trim}~\cite{drattrim}, which also emitted an LRAT proof that was re-verified by \texttt{lrat-check}~\cite{lrat}. The largest certificate, for $(13,7,5)$, has $446{,}686$ lemmas (a 37\,MB DRUP proof and an 82\,MB LRAT proof). For $n=13$ all instances were certified in the Q-only encoding. For $n=11$, $(11,6,4)$ was certified in the Q-only encoding ($35{,}510$ lemmas), and $(11,2..5,\cdot)$ in the encoding with Q and Lemma~T (Q-only is covered by the Lean proof below). A truncated proof is correctly rejected. On this route, the trusted components are the CNF encoder (a short Python program, cross-validated as in \cref{app:verification}) and the proof checkers.
\item \emph{$n=7$ without SAT.} A plain backtracking enumeration of all colourings, with Q checked as soon as its pairs are assigned, finds no ordinal counterexample for $D=2,3,4$. For $D=4$ it does so with Lemma~T switched off.
\item \emph{Lean~4.} The statements for $n=7$, $11$ and $13$ are formalised in Lean~4~\cite{lean4} with Mathlib~\cite{mathlib}; see below.
\end{itemize}

\subsubsection*{The Lean formalisation ($n=7$, $11$, $13$)}
The file \texttt{Erdos/Erdos/E132Convex.lean} proves, without \texttt{sorry},
\begin{center}
\texttt{theorem erdos132\_convex7 (p : Fin 7 $\to$ EuclideanSpace $\R$ (Fin 2)) (hp : StrictConvexPolygon p) : 3 $\le$ (rareDists p).card}
\end{center}
and the analogous \texttt{erdos132\_convex11} and \texttt{erdos132\_convex13} (file \texttt{E132Convex13.lean}) for \texttt{Fin 11} and \texttt{Fin 13}. Here \texttt{StrictConvexPolygon p} means that \emph{every} index-increasing triple $(p_i,p_j,p_k)$, $i<j<k$, is positively oriented; that is, the vertices are given as a strictly convex polygon listed counter-clockwise. \texttt{rareDists p} is the set of distances realised by pairs $i<j$ with multiplicity at most $n$. The proof formalises \cref{lem:Q} (via the intersection point of the diagonals and strict convexity of the Euclidean norm), the passage to ranks and Q-colourings, \cref{lem:R}, and the correctness of a bit-vector encoding of the finite statement ``every Q-colouring with colours below $D_{\max}$ has at least three classes of size in $[1,n]$''. This covers all $D\le D_{\max}$ at once, with unused colours allowed. The finite statement is closed by Lean's \texttt{bv\_decide}, which bit-blasts the formula with a verified bit-blaster, calls the bundled SAT solver CaDiCaL, and checks the resulting LRAT certificate with a checker that is proved correct in Lean.

\emph{Trust base.} \texttt{\#print\allowbreak{} axioms} reports, for each theorem, the standard axioms \texttt{propext}, \texttt{Classical.choice}, \texttt{Quot.sound}, plus one axiom generated by \texttt{bv\_decide} (\texttt{\dots\_native.\allowbreak bv\_decide.\allowbreak ax\_1\_14}). That axiom asserts that the verified LRAT checker \texttt{verifyBVExpr} returns \texttt{true} on the concrete formula and certificate. The value \texttt{true} is obtained by running compiled code, not by kernel reduction, so the trust base is comparable to that of \texttt{native\_decide}: it includes the Lean compiler and runtime. Neither our Python encoder, nor any external SAT solver, nor \texttt{drat-trim} is trusted on this route. An attempt to check the $n=7$ axiom statement by kernel reduction ran for more than 20 minutes without finishing and was stopped.

\emph{Unordered point sets.} The file \texttt{E132ConvexFinset.lean} proves (using only the three standard axioms) that a finite set $S$ in convex position, in the sense of Mathlib's \texttt{ConvexIndependent}, can be listed as a \texttt{StrictConvexPolygon}. The proof sorts $S\setminus\{p_0\}$ by orientation about a point $p_0\in S$; a negatively oriented increasing triple would put a point of $S$ inside a triangle spanned by others. It also proves that the rare distances of $S$ (unordered pairs) coincide with \texttt{rareDists p} for any such listing. This gives \texttt{erdos132\_convex\{7,11,13\}\_finset}: for \texttt{S : Finset} with \texttt{S.card} $=n$ and \texttt{ConvexIndependent}, $3\le$ \texttt{(rareDistsS S).card}.

\emph{The layer bound.} The results of \cref{sec:layers} are formalised in \texttt{E132Layer.lean} without \texttt{sorry}, and \texttt{\#print\allowbreak{} axioms} reports only \texttt{propext}, \texttt{Classical.choice} and \texttt{Quot.sound} (no \texttt{bv\_decide} or \texttt{native\_decide}). This covers the localisation of $\Dt$-pairs to $L_1\cup L_2$, the degree-$2$ property of $L_2$-vertices in the $2$-core, the bound $\deg\le4$ on hull vertices, Theorem~A, Corollary~A1 and its unconditional form. There $L_1$ is the set of extreme points of $X$, shown equal to the vertex set of $\operatorname{conv}X$, and the $2$-core is the largest subset with minimum degree at least $2$.

\emph{The main theorem.} Theorem~1$'$ (the bound $54/37$) is formalised in full in \texttt{E132Main*.lean}, \texttt{E132ConvexCurve.lean} and \texttt{E132Ve87.lean} (about 7{,}000 lines):
\begin{center}
\texttt{theorem erdos132\_main : $\exists$ C\textsubscript{0} : $\R$, $\forall$ X : Finset Pt, 2 $\le$ X.card $\to$ (min (mult X (dist2 X)) (mult X (minDist X)) : $\R$) $\le$ 54/37 * X.card + C\textsubscript{0}}
\end{center}
with \texttt{mult X d} the number of unordered pairs of distinct points of $X$ at distance $d$, and \texttt{dist2}, \texttt{minDist} the second-largest and smallest distance. There is no \texttt{sorry}, and \texttt{\#print\allowbreak{} axioms} reports only \texttt{propext}, \texttt{Classical.choice} and \texttt{Quot.sound}. The formal proof follows the proof of Theorem~1$'$ in \cref{sec:setup,sec:reductions,sec:bad,sec:regions,sec:degenerate,sec:assembly} (ownership/pairing in Region~II, and \cref{lem:lp2} in Region~III). Theorem~1 (the bound $15/11$, \cref{sec:upgrade,sec:loss1,sec:loss2,sec:proof15}) is not formalised. The formalisation also proves Vesztergombi's bound $\mult(\Dt)\le\frac32n$ and the Hopf--Pannwitz bound, so no result is used as a black box. In three places it takes a different route from the text: (a) Vesztergombi's bound is proved in the $2$-core (every degree-$4$ core vertex has two degree-$2$ core neighbours, then double counting), with no point deletion; (b) the bad-edge bounds of \cref{sec:bad} are obtained by counting over directions of outward normals instead of arc length, so \cref{lem:lenscorner} (lens arc length) is not needed; (c) the degenerate-regime incidence bound uses a local charging argument, and the bound on the number of mixed T-edges that follows \cref{lem:mixedT} is proved in a weaker form (at most $1200$ instead of $71$), which only changes the constant $C_0$.

\emph{Limitations of the formalisation.} (i) The \texttt{bv\_decide} axiom described above remains. (ii) Only the $n=7$ LRAT certificate (1.6\,MB) is stored and replayed with \texttt{bv\_check}; for $n=11,13$ the certificates (23\,MB and 97\,MB) are not stored, so compilation re-runs CaDiCaL (about 40\,s and 4\,min respectively). (iii) \Cref{lem:lenscorner} is not formalised; the formal proof of Theorem~1$'$ does not need it (see above). (iv) Theorem~1, and in particular \cref{sec:upgrade,sec:loss1,sec:loss2,sec:proof15}, is not formalised.

\begin{remark}\label{rem:ordinal}
\begin{enumerate}[leftmargin=2em,label=(\alph*)]
\item The method is tight at ``three'': $(11,6,3)$ is satisfiable (e.g.\ with class sizes $(12,9,12,12,6,4)$), as is $(9,5,2)$.
\item The ordinal statement fails for $n=5,6$: $(6,3,1)$ has an ordinal solution with profile $(9,3,3)$, so at $n=6$ geometry beyond \cref{lem:Q} is needed.
\item For $n=15$, all instances with $D\le7$ were found unsatisfiable by CaDiCaL (solver output only, not certified), but the hard instance $(15,8,6)$ was unresolved after more than 20 minutes. So $n=15$ is open.
\item Theorem~\ref{thm:1113} for $n=7$ also follows from the literature (\cref{prop:lit678}).
\end{enumerate}
\end{remark}

\begin{remark}[related work]
Independently, a Lean project on the Lean Pool~\cite{lyfar} proves that for point sets with coordinates in certain restricted fields (e.g.\ $\mathbb Q(\sqrt2)$), $n\le2^r$ where $r$ is the number of rare distances. This gives at least three rare distances for $n\ge8$ in those coordinate fields; it neither implies nor is implied by \cref{thm:convex-summary}. Other contributions on the forum~\cite{forum132} concern the first question of Problem \#132 (at least two rare distances) for small $n$ and general sets.
\end{remark}

\section{Open problems}\label{sec:open}

\subsection*{\texorpdfstring{Problem 1.6 of~\cite{CDL25}}{Problem 1.6 of CDL}}
We now know $9/7\le L\le15/11$, i.e.\ $1.2857\ldots\le L\le1.3636\ldots$. The lower bound construction~\cite{ienjoymath} is conjectured there to be optimal; it is consistent with our proof (\cref{rem:ienjoy}). In the proof of Theorem~1 the $\delta$-side bound $e(S)\le s+|Q|+O(1)$ loses nothing on that construction, and the whole remaining gap sits in the contribution $\sum_{y\in R}(\deg y+m_y)/2$ of the points outside $S$: we prove $\deg y+m_y\le10$ for all but $O(1)$ of them (\cref{thm:N1}), while the construction has $\deg y+m_y\le6$. If $\deg y+m_y\le6$ held for all but $o(n)$ points of $R$, the linear programme would give $9/7$ in Region~I, but only $4/3$ in Regions~II and~III, because of the term $|Q|$ (\cref{rem:lpgeneral}); so $9/7$ would also need a sharper bound on $\mult(\Dt)$ there. For $\tau<\sqrt3/2$ the bound $10$ of \cref{thm:N1} is sharp for local arguments (\cref{rem:sharpN1}): in the example given there, $y$ is the inner apex of an equilateral triangle on the $\delta$-edge $z_+z_-$ between two points that may lie in $D$, and this edge determines $y$. Whether $\Theta(n)$ points of weight $10$ can occur together with a large $\mult(\Dt)$ is open; charging such points to $\delta$-edges between points of $D$ is one possible approach.

\subsection*{Genericity}
One might hope that excluding four concyclic points forces $\mult(\Dt)\le n+O(1)$, which together with \cref{thm:layer} would answer the first question of Problem \#132 for such sets. This is false. For $n=4m\ge12$, start from $n/2$ points on a circle of radius $\frac12$ interleaved with $n/2$ points on a circle of radius $R$, where $R^2+R\cos(2\pi/n)=\frac34$, and perturb generically within the solution set of the distance equations. This gives convex sets with no four concyclic points and $\mult(\Dt)=\frac54n$. We have verified these examples to $50$ digits for $n=12,20,40$ (all equalities to $5\cdot10^{-51}$, with concyclicity determinants at least $10^{-8}$). A non-convex family with $\mult(\Dt)=\frac54n-\frac32$ was verified numerically for $n\le62$. In both families every distance other than $\Delta$ and $\Dt$ occurs once, so the first question of Problem \#132 still has a positive answer for them, but the second rare distance is not $\Dt$. An elementary observation, probably folklore, explains the role of concyclicity in the extremal configurations: if $|ab|=|cd|=\Delta$ and $|ac|=|bd|=\Dt$, then $a,b,c,d$ are concyclic. (The triangles $abc$ and $dcb$ are congruent. A half-turn would make $abdc$ a parallelogram with $|ad|^2+|bc|^2=2(\Delta^2+\Dt^2)>2\Delta^2$, which is impossible, so $abdc$ is an isosceles trapezoid.)
\begin{question}\label{q:genericity}
If no four points of $X$ are concyclic, is $\mult(\Dt)\le\frac54n+O(1)$?
\end{question}

\subsection*{Convex position}
\begin{conjecture}[Q-ORD]\label{conj:QORD}
For every $n\ge7$, every Q-colouring of the pairs of $\{0,\dots,n-1\}$ has at least three colour classes of size at most $n$.
\end{conjecture}
Q-ORD holds for $7\le n\le13$ (\cref{thm:1113}), and it would imply that every set of $n\ge7$ points in convex position has at least three rare distances. For $n=15$ only $(15,8,6)$ remains. The solving time grows roughly tenfold for each increase of $n$ by $2$, so a proof for all $n$ must be human. Deletion-minimal refutations are not local ($(7,4,2)$ already needs $25$ of the $35$ quadrilaterals), so such a proof will have to be global.

For odd $n=2m+1\ge7$ (no computation needed), \cref{lem:counting} and (Alt) show that a convex set with $k\le2$ must have $D=m+1$, $k=2$ and $m-1$ frequent distances. Moreover, writing the frequent multiplicities as $2m+2+e_d$ and the two rare ones as $a,b$, the identity $a+b+\sum_de_d=m+2$ holds. (Indeed $D=m$ gives $R_n$ with $k=m\ge3$, $D\ge m+2$ gives $k\ge3$, and the identity is $(m-1)(2m+2)+\sum e_d+a+b=m(2m+1)$.) So for odd $n\ge15$ the question in convex position is about this single kind of configuration. Beyond three, the second question of Problem \#132 asks whether $k\to\infty$; this is open even in convex position.

\appendix
\section{Computational verification}\label{app:verification}

All scripts are in the directory \texttt{scripts/} of the accompanying repository \url{https://github.com/g8r-b8/erdos132-lean}. None of the computations below is used as a proof step, except the SAT/LRAT certificates and the Lean proof for \cref{thm:1113}. The rest are independent checks of steps that are proved in the text. Tags \chk{Ck} in the text refer to sections of \texttt{writeup\_a2\_checks.py}.

\subsection*{Theorems 1 and 1\texorpdfstring{$'$}{'}}
The following checks concern the proof of Theorem~1$'$ (\cref{sec:setup,sec:reductions,sec:bad,sec:regions,sec:degenerate,sec:assembly}).
\begin{itemize}[leftmargin=2em]
\item \texttt{writeup\_a2\_checks.py}: 58 assertions in exact rationals and rational interval arithmetic, with one-sided rational bounds for $\sqrt{\cdot}$, $\pi$, $\sin$, $\cos$, $\arcsin$. Sections: C1 the LP certificates of \cref{lem:lp,lem:lp2}; C2 the identity of \cref{lem:identities}(c) and the relation between $\tau$ and $t$; C3 small-angle facts for $\Dt\ge10^4$; C4 the Region~I windows; C5 the Region~II constants ((R1), (R2), (R2-sym), mixed T-edges, rhombus threshold, pairing ratio); C6 the corner lemma; C7 the lens constants; C8 case (c2); C9 the bad-edge totals; C10 the degenerate-regime constants; C11 the small-case packing bound. All pass.
\item \texttt{ref2\_numeric.py}: an independent set of 40 exact assertions (second referee pass) covering the same inequalities, together with an adversarial floating-point search for (R1), (R2), (R2-sym). The search is reported, not asserted. Its maxima were $0.906$, $0.897$ and $0.914$, against the proved bounds $0.908$, $0.918$ and $0.948$.
\item \texttt{ref2\_ownership.py}: an exhaustive exact check of the ownership/pairing bookkeeping in \cref{sec:regions}, in the abstract model the argument uses: rows that are disjoint unions of paths and cycles of length $\ge3$, rungs forming a matching, (R2), (R2-sym), and no two consecutive rhombi. It covers all row structures up to isomorphism with at most $6$ vertices per row and all partial rung matchings. There were $2{,}365{,}426$ admissible configurations without rhombi, with maximum ratio exactly $5/4$, and $2{,}741{,}183$ with rhombi allowed, with maximum ratio exactly $4/3$. Mutation tests show that (R2-sym) and the rule excluding two consecutive rhombi are both needed.
\item \texttt{ref\_degen.py}, \texttt{degen\_check.py}, \texttt{ref2\_degen.py}: adversarial searches in the degenerate regime, including $5500$ configurations built so that three or more circles $\Cc(w,\Dt)$ pass through one vertex of $\partial K$. The maximum of $\mult(\Dt)-n$ found was $0$, against the proved bound $1$.
\item \texttt{ref\_checks.py}, \texttt{a2\_checks.py}: random tests of the lens lemma (about $7000$ lenses, worst ratio $0.66$) and of the exact single-centre strip in the near-resonant regime.
\end{itemize}
The whole proof of Theorem~1$'$ was also checked twice more, independently, in adversarial ``referee'' passes. Vesztergombi's theorem~\cite{Ve87} was checked in the original: its main theorem is $\mult(\Dt)\le\frac32n$ for all $n$, with no additive constant.

The following checks concern the upgrade to Theorem~1 (\cref{sec:upgrade,sec:loss1,sec:loss2,sec:proof15}).
\begin{itemize}[leftmargin=2em]
\item \texttt{r4\_regionII\_checks.py}: $17$ assertions in exact rational arithmetic for Region~II: the count of type-(i) vertices, the identity and threshold of \cref{lem:farpartner} (on $500$ exact rational instances), a sufficient condition for the inequalities $\eta,\eta^*<t$ in \cref{lem:doublerung} on all of Region~II (an alternative to the angle condition used in the text), and the $18/13$ linear programme of \cref{rem:lpgeneral}. It also runs adversarial floating-point searches against \cref{lem:doublerung}, which are reported, not asserted.
\item \texttt{r4\_regionIII\_checks.py}: $18$ assertions in exact rational arithmetic for Region~III: the constants of \cref{lem:slot} and of the injectivity of killed slots, the corner-lemma thresholds, the mixed T-edge count $<76$, and the constants of \cref{prop:tauexact} (the bound $t<0.8661$ at $\tau=\sqrt3/2$, the threshold of \cref{lem:farpartner}, the bound $\eta<0.0201$ in \cref{lem:onecircle}, the bound on $\arcsin$, and $t+1/(23\Dt)<3/\pi$ in \cref{lem:arith}). Three floating-point searches are reported: B1, the slot-killing distance (maximum $|y-z|=0.0691$, against the proved bound $0.0708$); B2, the step of \cref{lem:rigid} ($0$ violations of $v_{i+1}=v^*$ in $2\cdot10^5$ samples); B3, the nearest approach $\min|j\alpha_2-k\alpha_1|\,\Dt=1.7\cdot10^{-6}$ (nonzero, as \cref{lem:arith} requires).
\item \texttt{r4\_lp\_checks.py}: exact certificates for the linear programmes of \cref{lem:lp3} and \cref{rem:lpgeneral}, and a re-check of those of \cref{lem:lp,lem:lp2}.
\item \texttt{r4\_nb\_exact.py}: in exact rational and $60$-digit decimal arithmetic, the identity $c_D(0)=\tau$ and the monotonicity of $c_D$ in \cref{rem:sharpN1}, the linear programme values, the constants of \cref{lem:thin,lem:bigturn} with rigorous rational bounds for $\sin$, $\cos$ and $\tan$, and the example of \cref{rem:sharpN1}.
\item \texttt{r4\_nb\_referee\_checks.py}: an independent set of exact checks of \cref{thm:N1,thm:N2}: the linear programmes, the apex-arc bound $2+2/\sin(0.5^\circ)<231.19<232$, the fatness margin and $W<9$ in \cref{lem:thin}, the multiplicity $M\le468^2=219024$ and the exceptional count in \cref{lem:bigturn} ($<6.88\cdot10^7$).
\item \texttt{r4\_nb\_search.py}: an adversarial floating-point search over random local boundaries of $K$ with corners, testing the neighbour bounds of \cref{thm:N1}; $0$ violations. It is a screening test only.
\item \texttt{r4\_nb\_ring16\_check.py}: consistency with the construction of~\cite{ienjoymath} for $n=60,120,250$: every point of $R$ has $m_y=0$ and $\deg y+m_y\le6$ (floating point).
\end{itemize}
The new parts of the proof were checked in independent referee passes: Region~II (\cref{prop:regionII2}), one pass; Region~III with $\tau\ne\sqrt3/2$, one pass; Region~III with $\tau=\sqrt3/2$ (\cref{prop:tauexact}), two passes; \cref{thm:N1,thm:N2}, one pass.

\subsection*{\texorpdfstring{Section~\ref{sec:layers}}{Layer bound}}
\begin{itemize}[leftmargin=2em]
\item \texttt{va1b\_exact.py} (sympy, mpmath interval arithmetic): exact computation of the six Case~1 candidates of \cref{prop:core2} and their pairwise squared distances, with an exhaustive check that no compatible set of candidates covers all three $v_i$. It also gives a closed form and a rigorous interval check ($6000$ subintervals, $0$ failures) of the Case~2 branch $|r-v_2|=\Delta$.
\item \texttt{va1c\_sympy.py}: symbolic checks of the identity in \cref{lem:M}, of $\lambda(\lambda+2\cos\gamma)=0$, of the chord formula, and of the bisector direction in \cref{lem:F}.
\item \texttt{va1c\_local.py}: exhaustive search for the configuration of \cref{lem:F} with unit vectors from $\mathbb Z[\omega]$ (norm $\le147$, exact), $\mathbb Z[i]$ (norm $\le325$, exact) and $N$-th roots of unity ($5\le N\le60$): $0$ hits in about $3.2\cdot10^6$ configurations. A control run with weakened hypotheses does find hits.
\item \texttt{va1\_clique.py}, \texttt{va1\_cyclo.py}, \texttt{va1c\_random.py}: searches over maximal compatible sets in lattice and cyclotomic pools, plus greedy floating-point constructions. Every set was tested against \cref{lem:location}, \cref{prop:core2}, \cref{thm:deg4}, \cref{thm:layer} and \cref{cor:A1,cor:uncond}: more than $800{,}000$ sets in the first searches and about $81{,}000$ in the second, with $0$ violations.
\end{itemize}

\subsection*{\texorpdfstring{Section~\ref{sec:convex}}{Convex position}}
\begin{itemize}[leftmargin=2em]
\item \texttt{ver\_classification\_k.py}, \texttt{ver\_hexagons.py}: exact chord-index profiles for the classified extremal sets used in \cref{thm:even,thm:nine,prop:lit678}, and strict convexity and multiplicity profiles of the three $3$-distance hexagons.
\item \texttt{counting\_parity.py}: the closed forms of \cref{lem:counting} for $n<400$.
\item \texttt{b3\_sat.py}: the CNF encoder (direct and order encodings, with or without Lemma~T). \texttt{b3\_enum.py}: the standalone backtracking enumerator for $n=7$.
\item \texttt{b3\_validate.py}: encoder validation. Every $7$-subset of $R_N$, $7\le N\le16$ ($11{,}440$ sets), and every $11$-subset of $R_N$, $N\le15$, satisfies all constraints under exact chord-index ranks, as do $3000$ (for $n=7$) and $500$ (for $n=11$) random convex polygons. The $(7,3,0)$ enumeration returns exactly the $R_7$ colouring. Every satisfying assignment printed by the SAT route was re-checked against the independent constraint code of \texttt{b3\_enum.py}.
\item Four solvers (CaDiCaL, Glucose~4, MiniSat~2.2, Lingeling) agree on the $n=7$ and $n=11$ instances.
\item \texttt{b3\_drup\_check.py}: an independent forward RUP checker, used for the $n=7$, $n=9$ and $n=11$ DRUP proofs. Its self-test rejects an empty proof of a satisfiable instance, a proof with one negated lemma, and a truncated proof.
\item \texttt{b3\_drat\_trim.py}: regenerates each CNF and Glucose proof deterministically and runs \texttt{drat-trim} with LRAT output, followed by \texttt{lrat-check}. The CNF, DRUP and LRAT files are not stored in the repository because of their size.
\end{itemize}

\begin{table}[ht]
\centering
\renewcommand{\arraystretch}{1.15}
\begin{tabular}{@{}llrll@{}}
\toprule
instance $(n,D,F)$ & encoding & lemmas & \texttt{drat-trim} & \texttt{lrat-check}\\
\midrule
$(7,4,2)$ & Q & 714 & verified & verified\\
$(11,2,0)$, $(11,3,0)$, $(11,4,0)$ & Q+T & 3, 8, 56 & verified & verified\\
$(11,5,3)$ & Q+T & 524 & verified & verified\\
$(11,6,4)$ & Q & 35,510 & verified & verified\\
$(13,2,0)$, $(13,3,0)$, $(13,4,0)$, $(13,5,0)$ & Q & 2, 7, 82, 448 & verified & verified\\
$(13,6,4)$ & Q & 8,463 & verified & verified\\
$(13,7,5)$ & Q & 446,686 & verified & verified\\
\bottomrule
\end{tabular}
\caption{Certified refutations for \cref{thm:1113}.}
\end{table}

\emph{Lean.} The whole development (all \texttt{E132*.lean} files: the convex-position results, the layer bound, Vesztergombi's and Hopf--Pannwitz's bounds, and the main theorem) builds with \texttt{lake exe cache get} followed by \texttt{lake build}; on a laptop this takes about 6 minutes after the Mathlib cache is downloaded, of which about 4 minutes is \texttt{bv\_decide} for $n=13$. The toolchain is Lean \texttt{v4.35.0-rc2} with Mathlib.

\bibliographystyle{amsplain}
\bibliography{refs}

\end{document}
